\documentclass[11pt,a4paper]{article}
\usepackage[T1]{fontenc}
\usepackage[utf8]{inputenc}
\usepackage{lmodern}
\usepackage[a4paper,margin=25mm,headheight=15pt]{geometry}
\usepackage{amsmath,amssymb,amsthm,mathtools,aliascnt}
\usepackage{microtype,booktabs,tabularx,array,longtable}
\usepackage[dvipsnames]{xcolor}
\usepackage{enumitem,fancyhdr,titlesec}
\usepackage[colorlinks=true,linkcolor=MidnightBlue,citecolor=MidnightBlue,urlcolor=MidnightBlue]{hyperref}
\usepackage[nameinlink,noabbrev]{cleveref}
\definecolor{Ink}{HTML}{243747}
\definecolor{Accent}{HTML}{566348}
\titleformat{\section}{\Large\bfseries\color{Ink}}{\thesection}{.65em}{}
\titleformat{\subsection}{\large\bfseries\color{Accent}}{\thesubsection}{.65em}{}
\pagestyle{fancy}\fancyhf{}
\fancyhead[L]{\small Surcomplex spectral theory}
\fancyhead[R]{\small Singular values and scales}
\fancyfoot[C]{\thepage}
\setlength{\parskip}{.28em}
\setlength{\parindent}{1.1em}
\setlist{itemsep=.15em,topsep=.35em,leftmargin=2em}
\setcounter{tocdepth}{2}
\emergencystretch=2em
\numberwithin{equation}{section}
\newtheorem{theorem}{Theorem}[section]
\newaliascnt{proposition}{theorem}
\newtheorem{proposition}[proposition]{Proposition}
\aliascntresetthe{proposition}
\newaliascnt{lemma}{theorem}
\newtheorem{lemma}[lemma]{Lemma}
\aliascntresetthe{lemma}
\newaliascnt{corollary}{theorem}
\newtheorem{corollary}[corollary]{Corollary}
\aliascntresetthe{corollary}
\theoremstyle{definition}
\newaliascnt{definition}{theorem}
\newtheorem{definition}[definition]{Definition}
\aliascntresetthe{definition}
\newaliascnt{example}{theorem}
\newtheorem{example}[example]{Example}
\aliascntresetthe{example}
\theoremstyle{remark}
\newaliascnt{remark}{theorem}
\newtheorem{remark}[remark]{Remark}
\aliascntresetthe{remark}
\newaliascnt{warning}{theorem}
\newtheorem{warning}[warning]{Warning}
\aliascntresetthe{warning}
\crefname{theorem}{Theorem}{Theorems}
\crefname{proposition}{Proposition}{Propositions}
\crefname{lemma}{Lemma}{Lemmas}
\crefname{corollary}{Corollary}{Corollaries}
\crefname{definition}{Definition}{Definitions}
\crefname{example}{Example}{Examples}
\crefname{remark}{Remark}{Remarks}
\crefname{warning}{Warning}{Warnings}
\crefname{section}{Section}{Sections}
\newcommand{\No}{\mathbf{No}}
\newcommand{\SC}{\mathbf{No}[i]}
\newcommand{\R}{\mathbb R}
\newcommand{\C}{\mathbb C}
\newcommand{\Q}{\mathbb Q}
\newcommand{\N}{\mathbb N}
\newcommand{\Z}{\mathbb Z}
\newcommand{\OO}{\mathcal O}
\newcommand{\mm}{\mathfrak m}
\newcommand{\supp}{\operatorname{supp}}
\newcommand{\st}{\operatorname{st}}
\newcommand{\lc}{\operatorname{lc}}
\newcommand{\rank}{\operatorname{rank}}
\newcommand{\tr}{\operatorname{tr}}
\newcommand{\diag}{\operatorname{diag}}
\newcommand{\spec}{\operatorname{spec}}
\newcommand{\im}{\operatorname{im}}
\newcommand{\GL}{\operatorname{GL}}
\newcommand{\Span}{\operatorname{span}}
\newcommand{\Hom}{\operatorname{Hom}}
\newcommand{\Gal}{\operatorname{Gal}}
\newcommand{\Res}{\operatorname{Res}}
\newcommand{\lcm}{\operatorname{lcm}}
\newcommand{\op}{\mathrm{op}}
\newcommand{\Fr}{\mathrm F}
\newcommand{\abs}[1]{\left|#1\right|}
\newcommand{\norm}[1]{\left\lVert#1\right\rVert}
\newcommand{\inn}[2]{\left\langle#1,#2\right\rangle}
\newcommand{\repocommit}{e260237db9b71da8b74a0c13c8e6355119091100}
\newcolumntype{L}{>{\raggedright\arraybackslash}X}
\newcolumntype{P}[1]{>{\raggedright\arraybackslash}p{#1}}
\hypersetup{pdftitle={Surcomplex Spectral Theory: Singular Values, Infinitesimal Rank, and Multiscale Stability},pdfauthor={Research exposition prepared for Vladimir Reshetnikov},pdfsubject={Finite-dimensional linear algebra over real closed fields and surreal Hahn workspaces}}
\begin{document}
\hypersetup{pageanchor=false}
\begin{titlepage}
\centering
\vspace*{12mm}
{\large\scshape A missing chapter for the Surreal documentation\par}
\vspace{12mm}
{\Huge\bfseries\color{Ink} Surcomplex\par Spectral Theory\par}
\vspace{7mm}
{\LARGE Singular Values, Infinitesimal Rank,\par and Multiscale Stability\par}
\vspace{9mm}
{\large September 21, 2026\par}
\vspace{11mm}
\begin{minipage}{.94\textwidth}
\begin{abstract}
Finite matrices over the surcomplex numbers admit much of the familiar
spectral theory, but its justification is algebraic rather than compactness
based, and its interpretation changes when singular values occupy different
infinitesimal scales. We develop Hermitian and normal diagonalization,
positive square roots, singular value and polar decompositions, variational
principles, least squares, low-rank approximation, conditioning, and
spectral-subspace perturbation over a real closed field and its
complexification.

The central synthesis identifies the valuations of singular values with
successive differences of the minimum valuations of minors. A positive
Cauchy--Binet identity then recovers both these scales and the leading
singular coefficients from the characteristic coefficients of the Gram
matrix. This gives an intrinsic rank filtration, an exact root-free scale
algorithm, and precision certificates valid for arbitrary-rank Hahn value
groups. Worked examples separate exact rank from standard-part rank,
small perturbations from gap-relative perturbations, and algebraic inverses
from Hahn-summable inverse series. Nonnormal pseudospectra, regularization,
and a finite tensor application complete the picture.

Divisibility of the value group is then removed. A finite splitting tree,
licensed by strong Hahn summability rather than by any convergence argument,
produces Hermitian and normal diagonalizations and singular-value
decompositions whose every exponent lies in the subgroup generated by the
input exponents. Against this, positive matrix roots genuinely cost new
scales, and the cost is exactly the determinantal one: for positive
semidefinite $A$ of rank $r$ the principal $q$th root generates the Hahn
field on $\Gamma+\sum_k\Z\Delta_k(A)/q$, a finite abelian totally ramified
extension of degree the order of the subgroup the classes $\Delta_k$ span in
$\Gamma/q\Gamma$; it is trivial exactly when every $\Delta_k$ lies in
$q\Gamma$, which is also the exact criterion for a Gram factorization
$A=C^*C$ over the base field. The single scalar $\tr(A^{1/q})$ is always a
primitive generator of that whole extension, by a positivity mechanism
rather than by genericity.
\end{abstract}
\end{minipage}
\vfill
\begin{minipage}{.94\textwidth}\small
\textbf{Status and scope.} A self-contained research exposition relative to
standard surreal normal-form and Hahn-field foundations, assembled from two
separately delivered manuscripts (see \cref{sec:gap}). Classical results
are identified as such; no priority claim or solution of a named open problem
is made. Older power-series spectral theorems, Hahn support lemmas and
Kummer theory are acknowledged as predecessors, and the arbitrary-rank
determinantal and primitive-trace package of
\cref{spec:sec:ramification,spec:sec:trace} is offered as a candidate
original contribution whose priority is not certified. All matrix dimensions
are ordinary finite integers. General proofs are not proof-assistant
verified; no Lean formalization of any result here is included or claimed.
The accompanying exact symbolic checks exercise finite examples, not
arbitrary Hahn supports or the general theorems.
\end{minipage}
\end{titlepage}
\hypersetup{pageanchor=true}
\pagenumbering{roman}
\tableofcontents
\clearpage
\pagenumbering{arabic}

\section{The documentation gap and the mathematical agenda}\label{sec:gap}

\subsection{What was inspected}
This article was chosen after inspecting the documentation catalogue and
relevant inventories in the repository
\href{https://github.com/VladimirReshetnikov/Surreal}{\texttt{VladimirReshetnikov/Surreal}},
at the fixed revision
\begin{center}\small\texttt{\repocommit}.\end{center}
The catalogue lists fifteen packages in surreal mathematics, surcomplex
analysis and geometry, and foundations and computation~\cite{RepoMap}.
The detailed polynomial-algebra inventory and its opening source sections
were especially relevant~\cite{RepoPoly}; the foundations inventory was
also inspected~\cite{RepoFound}. This is a targeted coverage audit, not a
claim to have checked every line of all archived source manuscripts.

Matrix theory is not wholly absent. The polynomial report uses an algebraic
Schur inequality and a differentiating compression in its treatment of
Schoenberg's inequality, and it uses signatures of Hermite forms to count
real roots. Finite residue pairings and multiplication matrices also appear
in the surrounding algebra. These are substantial ingredients, but they are
organized around polynomial roots and residues, not around matrices as the
objects of study. The inspected inventories do not offer a systematic
singular-value, conditioning, or scale-resolved rank theory.

\begin{table}[ht]
\centering\small
\begin{tabularx}{\textwidth}{@{}p{.24\textwidth}LL@{}}
\toprule
Existing strand & What it supplies & What this chapter adds\\
\midrule
Polynomial algebra & Root geometry, Schur-type arguments, signatures,
Newton profiles & Unitary geometry, SVD, positive Gram matrices, determinantal
singular scales\\
Analytic geometry & Finite algebras and deformation questions & Gap-sensitive
spectral projections and Schur-complement effective matrices\\
Foundations & Set-sized Hahn workspaces; limits on transfer & Explicit finite
algebraic proofs without compactness or a class-sized choice principle\\
Computer algebra & Representation and precision distinctions & Root-free
scale elimination and certificates for resolvable rank\\
\bottomrule
\end{tabularx}
\caption{The gap is an organizing theory, not the absence of every matrix lemma.}
\end{table}

\subsection{Why this is more than replacing scalars by surreals}
Over the ordinary reals, a positive singular value is often tacitly regarded
as a non-negligible number. Over the surreals, a matrix can have singular
values at scales $1$, $\omega^{-2}$, and $\omega^{-\omega}$. It is then
simultaneously invertible, singular after taking standard parts, and
indistinguishable from a lower-rank matrix at many prescribed precisions.
An infinitesimal perturbation can rotate an eigenspace through an ordinary
angle when its spectral gap is infinitesimal too. Conversely, an infinite
entry need not imply an unstable direction if every singular value has the
same scale.

The issue is therefore not just whether a decomposition exists. We need to
know what its invariants measure, which conclusions survive reduction to an
ordinary complex matrix, and which operations can be certified without
computing an entire Hahn expansion. The key bridge is
\begin{equation}\label{eq:introbridge}
 \min_{|I|=|J|=k}v\bigl(\det A_{I,J}\bigr)
   =\sum_{j=1}^k v\bigl(\sigma_j(A)\bigr),
 \qquad 1\le k\le\rank A.
\end{equation}
Here the singular values are ordered by their actual positive magnitudes,
whereas the valuations are ordered in the opposite sense. This identity
turns analytic-looking size information into finite determinantal algebra.

\subsection{Attribution and reading route}
The spectral theorem over real closed fields is classical, not a novel
surreal phenomenon; a recent explicit use is Wang~\cite{Wang}. The
variational and perturbation principles have their classical matrix-analysis
counterparts, for example in Tao's lecture notes~\cite{Tao}, the original
low-rank approximation paper of Eckart and Young~\cite{EY}, and the
Davis--Kahan literature~\cite{DK,YWS}. The pseudoinverse is the classical
Moore--Penrose construction~\cite{Penrose}.

The contribution here is an integrated account with proofs in the exact
field setting, the determinantal and Gram-coefficient description of
surreal scales, and explicit boundaries for reduction, computation, and
summability. We do not claim that these consequences of existing algebra
first occur here. Readers already comfortable with real-closed-field linear
algebra can start at \cref{sec:valuation}; readers interested primarily in
computation can proceed from there to \cref{sec:algorithm,sec:conditioning}.

\subsection{Provenance and the second half of the article}
\label{spec:sec:provenance}
This article is assembled from two separately delivered manuscripts, both
dated 21 September 2026 and both kept verbatim in \texttt{sources/}.
Every section except
\cref{spec:sec:descent,spec:sec:ramification,spec:sec:trace} comes from
\texttt{01-determinantal-singular-value-scales.tex}, which worked over a real
closed field and, for the scale-sensitive results, in a \emph{divisible} Hahn
workspace. Those three sections come from
\texttt{02-determinantal-ramification-descent.tex}, which removes that
divisibility hypothesis and measures exactly what it was buying.

The second manuscript cites the first by name and declines to reclaim its
content: the SVD, the determinantal singular scales of \cref{thm:scales} and
the positive Cauchy--Binet identity of \cref{thm:gram} are stated here once
only, in the first manuscript's form. What the second adds is the removal of
divisibility from the decompositions, and then an exact field-of-definition
and degree calculation for positive matrix roots, in which the very invariant
$\Delta_k$ that measures singular scales turns out to control a finite
radical extension.

The two manuscripts audited the repository at different pinned revisions,
\begin{center}\small
\texttt{\repocommit}\\
\nolinkurl{aa846271b4dcae2c055b216126a87210292ec19b}
\end{center}
the first and the second respectively. Both declared their audits targeted
rather than line-by-line, with absence from a search result never used as
proof of absence from the collection.

That second half has genuine predecessors, which it concedes rather than
displaces. Adkins proved a spectral theorem for normal matrices over
hermitian discrete valuation rings, including one-variable formal and
convergent power series~\cite{Adkins}; Keller and Ochsenius proved orthogonal
diagonalization over iterated real Laurent-series fields and a particular
infinite-rank Hahn field by recursive block splitting~\cite{KO}. So neither
``a spectral theorem beyond $\R$'' nor ``no fractional exponents in a
Laurent-series eigenbasis'' is claimed as new. Parusi\'nski and Rond prove
multiparameter formal spectral and singular-value results under
normal-crossing discriminant hypotheses~\cite{PR}, and the 2026 preprint of
Dai, Liang, Lu and Zhi gives a splitting-and-projector criterion over
multivariable formal power-series rings~\cite{DLLZ}. A Hahn field is a
different ring: a totally ordered exponent group permits monomial division
that a multivariable power-series ring need not permit, and we do not claim
to remove the obstructions in that latter category
(\cref{spec:warn:crossing}). The support foundations are standard
Hahn--Neumann and word-order ideas~\cite{Neumann,Higman}, reproved here for
self-containment, and the finite extension calculation is an elementary
monomial form of Kummer theory.

\subsection{Which spectrum this is}
\label{spec:sec:spectrumword}
The collection uses the word \emph{spectrum} for several inequivalent things,
so this article fixes its own sense once. Here $\spec H$ is the algebraic
\emph{eigenvalue list} of a finite matrix over $K$, together with the
valuations of those eigenvalues and of the singular values. It is not the
differential \emph{phase spectrum} of
\texttt{docs/surcomplex/differential-equations/}, which is a multiset of
purely infinite surreal parts of primitives of eigenvalues, and whose members
are not numbers of the same kind as ours and are not finite angles. It is not
the \emph{prime spectrum} of a ring, the sense in which
\texttt{docs/surcomplex/analytic-geometry/} discusses maximal and prime
ideals. And the visible-rank levels $r_{<\theta},r_{\le\theta}$ of
\cref{sec:rank}, although they form a list of scales, are not a spectrum in
any of these senses.

\section{Fields, workspaces, and conventions}\label{sec:fields}

\subsection{The algebraic setting}
Until \cref{sec:valuation}, let $F$ be a set-sized real closed ordered field
and put $K=F(i)$, where $i^2=-1$. We use the standard equivalent
characterization that $F$ is ordered and $F(i)$ is algebraically closed.
Conjugation is $\overline{x+iy}=x-iy$, and
\[
 |z|=\sqrt{z\bar z}=\sqrt{x^2+y^2}\in F_{\ge0}.
\]
Every positive element of $F$ has a unique positive square root. No
Archimedean property or completeness assumption is imposed.

For a matrix $A$, write $A^*=\bar A^{\mathsf T}$. A square matrix is
\emph{Hermitian} when $A^*=A$, \emph{normal} when $A^*A=AA^*$, and
\emph{unitary} when $A^*A=AA^*=I$. A matrix with real entries is treated
inside $F$, with transpose in place of conjugate transpose. The integers
$m,n,k$ used as dimensions and ranks are always ordinary finite integers.
Ambient row and column sizes are positive; zero-dimensional invariant
subspaces and empty products are allowed.
Unless explicitly stated otherwise, vector spaces and subspaces are over
$K$.

We use $\infty$ only as an added valuation symbol, with
$v(0)=\infty$. It is not a surreal number. Expressions such as
$\infty-\infty$ will never be used.

\subsection{Hahn workspaces inside the surcomplex numbers}
For the scale-sensitive results, choose a set-sized additive
subgroup $\Gamma\subseteq\No$, with its inherited order, and set
\begin{equation}\label{eq:workspace}
 F_\Gamma=\R((t^\Gamma)),\qquad K_\Gamma=\C((t^\Gamma)),
 \qquad t^\gamma:=\omega^{-\gamma}.
\end{equation}
The last expression uses Conway's monomial map. It is not a convention
about the ordinary real exponential. An element of $K_\Gamma$ is a formal
sum
\[
 a=\sum_{\gamma\in S}a_\gamma t^\gamma,
 \qquad a_\gamma\in\C,
\]
where $S\subseteq\Gamma$ is a well-ordered set in increasing exponent
order. Zero coefficients may be omitted. The field operations are Hahn
addition and convolution. On $F_\Gamma$, the sign of a nonzero series is the
sign of its coefficient at the least exponent.

No divisibility, finite-rank, countability, or Archimedean assumption is
imposed on $\Gamma$.

We use two foundational facts. Conway normal forms identify these series
with surcomplex numbers, and the Hahn field with divisible value group and
algebraically closed coefficient field of characteristic zero is
algebraically closed. Gonshor's Chapter~5 gives the normal-form and support
foundations~\cite{Gonshor}; Poonen's Corollary~4 gives the relevant
algebraic-closedness statement for Mal'cev--Neumann fields~\cite{Poonen}.
Thus \emph{when $\Gamma$ is divisible}, $K_\Gamma=F_\Gamma(i)$ is
algebraically closed and $F_\Gamma$ is real closed, so every result proved
over a real closed $F$ applies verbatim inside the workspace. These standard
foundational theorems are inputs, not reproved here.

\begin{remark}[What still needs divisibility]\label{spec:rem:divisibility}
For a general $\Gamma$ the field $F_\Gamma$ is \emph{not} real closed:
if $\gamma\notin2\Gamma$ then the positive element $t^\gamma$ has no square
root in $F_\Gamma$, because the valuation of a square lies in $2\Gamma$.
Divisibility is nevertheless needed only where a root of an
\emph{arbitrary} positive element, or an algebraic closure, is actually
used: the Schur triangularization of a non-normal matrix in
\cref{prop:normal}, the positive semidefinite square root of an arbitrary
positive semidefinite matrix and the Cholesky factorization in
\cref{prop:positive} with the generalized Hermitian eigenproblem built on
it, and consequently the functional calculus and pseudospectral statements
that invoke an eigenvalue of an arbitrary square matrix --- all of which
live in the divisible hull $\Gamma\otimes_\Z\Q$, at a cost computed exactly
in \cref{spec:thm:ramification}. Everything else in this article, including
the Hermitian and normal spectral theorems, the SVD, the polar
decomposition, every norm and variational principle, the determinantal
scales of \cref{thm:scales}, the positive Cauchy--Binet identity of
\cref{thm:gram}, the rank filtration, the scale algorithm, the perturbation
certificates, and the matrix series of \cref{thm:series}, holds over an
arbitrary set-sized ordered abelian $\Gamma$, by the support-certified
constructions of \cref{spec:sec:descent}.
\end{remark}

\begin{proposition}[Finite-data localization]\label{prop:localization}
Every finite collection of surcomplex matrices has all its entries in one
workspace \eqref{eq:workspace}. All the decompositions established below can
be performed in that workspace.
\end{proposition}
\begin{proof}
Take the union of the exponent supports in the real and imaginary normal
forms of every entry. There are finitely many entries and every support is
a set, so their union is a set. Its rational linear span is a divisible
set-sized subgroup of $\No$. Reversing exponents to follow the convention
$t^\gamma=\omega^{-\gamma}$ puts all entries in its Hahn field.
Algebraic closedness supplies all roots needed in the spectral proofs;
positive square roots and finite field operations remain in the same
workspace. Every basis extension used below has finite length.
\end{proof}

The same argument works for a set of scalar inputs, but we only need finite
matrices here. No quantification over a class-sized set of all bases is
required. A displayed theorem about matrices over $\SC$ means that the
finite data are localized and the theorem is applied there. Its explicit
algebraic witnesses are still witnesses in any larger workspace.

\Cref{prop:localization} deliberately passes to a rational span, which is
the smallest \emph{divisible} workspace containing the data. There is a
smaller and more informative choice.

\begin{proposition}[Exact-support localization]\label{spec:prop:exactlocalization}
Let $A$ be a finite surcomplex matrix. Take the union of the exponent
supports of the real and imaginary Conway normal forms of all its entries,
negate the exponents to follow the convention $t^\gamma=\omega^{-\gamma}$,
and let $H\subseteq\No$ be the additive subgroup they generate. Then $H$ is
a set, every entry of $A$ lies in $K_H$, and $H$ is contained in every
subgroup of $\No$ whose Hahn field contains all the entries.
\end{proposition}
\begin{proof}
A subgroup generated by a set consists of the finite integer combinations of
its generators, so it is a set. The remaining assertions are immediate from
the definition of the support of a normal form.
\end{proof}

The group $H$ need not be divisible, and that is the point: the results of
\cref{spec:sec:descent} take place inside $K_H$ itself, while
\cref{spec:sec:ramification,spec:sec:trace} measure exactly how far a
positive matrix root has to leave it. No proper-class Hahn sum is used
anywhere in this article.

\subsection{Order, valuation, and topology are separate structures}
In a Hahn workspace the ordered modulus and the valuation are not the same
thing. For $\rho>0$, both $t^\rho$ and $2t^\rho$ have valuation $\rho$,
although the latter is outside the ordered disk $|z|\le t^\rho$.
We will use ordered inequalities for optimization and conditioning, and
valuation inequalities for relative scale and support precision.

The principal proofs in this article use no infinite limiting process. A
unit sphere is not assumed compact, a supremum is not assumed to exist, and
no spectral integral is invoked. In \cref{sec:functions}, where an infinite
matrix series is used, its Hahn support is checked explicitly. Internal
valuation convergence in a chosen Hahn field and the fine order/modulus
topology inherited from the whole surreal class must not be conflated.

\section{Hermitian geometry without compactness}\label{sec:geometry}

For $x,y\in K^n$, define
\[
 \inn xy=x^*y=\sum_{j=1}^n\bar x_jy_j,
 \qquad \norm x=\sqrt{x^*x}.
\]
This convention is conjugate-linear in the first variable. Since $F$ is
ordered, a sum of squares is zero only if every term is zero. Consequently
$x^*x>0$ for $x\ne0$.

\begin{lemma}[Cauchy--Schwarz and orthogonal decomposition]\label{lem:CS}
For $x,y\in K^n$,
\[
 |x^*y|\le\norm x\norm y.
\]
If $x\ne0$, the vector $y$ decomposes uniquely as
\[
 y=x\frac{x^*y}{x^*x}+y_\perp,\qquad x^*y_\perp=0.
\]
Every finite-dimensional subspace admits an orthonormal basis and has an
orthogonal complement.
\end{lemma}
\begin{proof}
The asserted decomposition follows by substitution. Its squared norm is
\[
 y^*y=\frac{|x^*y|^2}{x^*x}+y_\perp^*y_\perp
       \ge\frac{|x^*y|^2}{x^*x},
\]
which proves Cauchy--Schwarz; the case $x=0$ is immediate. Gram--Schmidt now
works on any finite linearly independent list: subtract the finitely many
orthogonal components already constructed and divide the remaining nonzero
vector by its positive norm. Extend a basis of the subspace to a basis of
$K^n$ and continue this process. The remaining orthonormal vectors span its
orthogonal complement.
\end{proof}

This proof still works when a norm is infinitesimal. Dividing by that norm
may produce infinite coordinates at an intermediate step, but it is a legal
field operation; the final unit vector has all coordinate moduli at most
one. It is invalid to delete a nonzero vector merely because its standard
part vanishes.

The triangle inequality follows from Cauchy--Schwarz:
\[
 \norm{x+y}^2
 \le\norm x^2+2\norm x\norm y+\norm y^2.
\]
Finite orthogonal projections therefore exist by an explicit basis formula,
not by an infinite-dimensional projection theorem. If $u_1,\ldots,u_k$ is
an orthonormal basis of $W$, then
\begin{equation}\label{eq:proj}
 P_W=\sum_{j=1}^ku_ju_j^*,\qquad
 P_W^2=P_W=P_W^*,\qquad K^n=W\oplus W^\perp.
\end{equation}
The projector is independent of the chosen orthonormal basis because it
acts as the identity on $W$ and as zero on $W^\perp$.

\section{Spectral, positive, and singular-value decompositions}\label{sec:spectral}

\subsection{The finite spectral theorem}
\begin{theorem}[Hermitian spectral theorem]\label{thm:spectral}
For every Hermitian $H\in K^{n\times n}$, there are a unitary $U$ and
$\lambda_1,\ldots,\lambda_n\in F$ such that
\[
 H=U\diag(\lambda_1,\ldots,\lambda_n)U^*.
\]
For real symmetric $H$, $U$ can be chosen orthogonal with entries in $F$.
\end{theorem}
\begin{proof}
The characteristic polynomial splits in $K$, so there is an eigenpair
$Hx=\lambda x$ with $x\ne0$. The scalar $x^*Hx$ is fixed by conjugation,
hence lies in $F$, and
\[
 \lambda=\frac{x^*Hx}{x^*x}\in F.
\]
Normalize $x$. If $y\perp x$, then
$x^*Hy=(Hx)^*y=\lambda x^*y=0$, so $x^\perp$ is invariant under $H$.
The restriction to that complement is Hermitian. Induction on $n$ completes
the construction. For a real symmetric matrix, once $\lambda\in F$ is
known, singularity of $H-\lambda I$ over $F$ supplies a nonzero kernel
vector in $F^n$. The same induction can therefore be carried out over $F$.
\end{proof}

\begin{proposition}[Schur form and normal matrices]\label{prop:normal}
Every $A\in K^{n\times n}$ is unitarily similar to an upper triangular
matrix. It is unitarily diagonalizable if and only if it is normal.
\end{proposition}
\begin{proof}
Choose and normalize an eigenvector and extend it to an orthonormal basis.
The first column below the diagonal then vanishes. Apply induction to the
remaining square block to obtain upper triangular Schur form $T$.
If $T$ is normal, comparison of the $(1,1)$ entries of $TT^*$ and $T^*T$
gives
\[
 |t_{11}|^2+\sum_{j>1}|t_{1j}|^2=|t_{11}|^2.
\]
Positivity forces the rest of the first row to vanish. Induction on the
remaining normal block makes $T$ diagonal. Conversely a diagonal matrix is
normal, and unitary conjugation preserves normality.
\end{proof}

For distinct eigenvalues $\lambda$ of a Hermitian matrix, its spectral
projectors are finite polynomials:
\begin{equation}\label{eq:spectralprojector}
 P_\lambda=\prod_{\substack{\mu\in\spec H\\\mu\ne\lambda}}
       \frac{H-\mu I}{\lambda-\mu},
 \qquad H=\sum_{\lambda\in\spec H}\lambda P_\lambda.
\end{equation}
Evaluation on the diagonal form proves the formula. The eigenspace
projector is canonical; a basis within a repeated eigenspace is not.

\subsection{Positivity, square roots, and congruence}
A Hermitian matrix is positive semidefinite, written $H\succeq0$, when
$x^*Hx\ge0$ for every $x$, and positive definite, written $H\succ0$, when
this is strict for $x\ne0$. These are equivalent respectively to all
eigenvalues being nonnegative or positive. In particular, positive definite
does not mean bounded below by a positive \emph{ordinary real} multiple of
$I$.

\begin{proposition}[Positive square root and inertia]\label{prop:positive}
A positive semidefinite matrix has a unique positive semidefinite square
root. The numbers of positive, negative, and zero eigenvalues of a
Hermitian matrix are unchanged by an invertible congruence $H\mapsto S^*HS$.
\end{proposition}
\begin{proof}
Diagonalize $H$ and replace each eigenvalue by its nonnegative square root.
For uniqueness, a positive semidefinite $R$ with $R^2=H$ commutes with $H$,
so it preserves each eigenspace of $H$. On the eigenspace with eigenvalue
$\lambda$, diagonalization of $R$ gives only nonnegative roots of
$r^2=\lambda$, hence $R=\sqrt\lambda I$ there.

If $p$ is the number of positive eigenvalues of $H$, its positive eigenspace
is a $p$-dimensional subspace on which the quadratic form is positive
 definite. Any subspace of dimension greater than $p$ meets the span of the
nonpositive eigenvectors nontrivially, so cannot have that property. Thus
$p$ is an intrinsic maximum dimension, preserved by invertible congruence.
Apply the same argument to $-H$ for the negative count; subtract both counts
from $n$ for the nullity.
\end{proof}

The usual Cholesky factorization also survives. If
$H=\left(\begin{smallmatrix}a&b^*\\b&C\end{smallmatrix}\right)\succ0$,
then $a>0$ and $C-bb^*/a\succ0$: evaluate the form at
$(-a^{-1}b^*y,y)$. Inductively factor that Schur complement as $LL^*$.
Then
\[
 H=\begin{pmatrix}\sqrt a&0\\b/\sqrt a&L\end{pmatrix}
   \begin{pmatrix}\sqrt a&b^*/\sqrt a\\0&L^*\end{pmatrix}.
\]
This is a finite algorithm, not a limiting process.

\subsection{Singular values and the polar decomposition}
\begin{theorem}[Singular value decomposition]\label{thm:svd}
Let $A\in K^{m\times n}$ and $p=\min(m,n)$. There are unitaries $U,V$ and
a rectangular diagonal matrix $\Sigma$ such that
\[
 A=U\Sigma V^*,\qquad
 \sigma_1\ge\cdots\ge\sigma_r>0,
 \quad \sigma_{r+1}=\cdots=\sigma_p=0,
\]
where $r=\rank A$ and the diagonal entries of $\Sigma$ are the $\sigma_j$.
The positive $\sigma_j$ are uniquely determined with multiplicity.
\end{theorem}
\begin{proof}
The matrix $A^*A$ is positive semidefinite. Choose an orthonormal eigenbasis
$v_1,\ldots,v_n$ and write its positive eigenvalues as
$\sigma_1^2,\ldots,\sigma_r^2$. For $j\le r$, put
$u_j=Av_j/\sigma_j$. Then
\[
 u_j^*u_k=(\sigma_j\sigma_k)^{-1}v_j^*A^*Av_k=\delta_{jk}.
\]
When the eigenvalue is zero, $\norm{Av_j}^2=0$ implies $Av_j=0$.
Extend the $u_j$ to an orthonormal basis of $K^m$. These bases give the
claimed factorization, with precisely $r$ nonzero diagonal entries. The
characteristic polynomial of $A^*A$ determines the squared singular values,
and the positive square root determines the singular values themselves.
\end{proof}

When convenient we extend the list by zeros up to length $n$, so that
$\sigma_j^2$ lists all eigenvalues of $A^*A$. In particular, formulas using
input subspaces of dimension $k$ remain meaningful for $k>p$.

Set $|A|=(A^*A)^{1/2}$. The SVD gives the polar factorization
\begin{equation}\label{eq:polar}
 A=W|A|,\qquad
 W=\sum_{j=1}^ru_jv_j^*.
\end{equation}
Here $W$ is the unique map that is an isometry from $(\ker A)^\perp$ to
$\im A$, is zero on $\ker A$, and satisfies \eqref{eq:polar}. On that
complement $|A|$ is invertible, so the uniqueness follows directly. For a
square invertible matrix, $W$ is unitary. For a singular square matrix,
unitary extensions of $W$ exist but are generally not unique. Confusing the
canonical partial isometry with one of those extensions loses information.

\subsection{A generalized Hermitian eigenproblem}
If $G\succ0$ and $H=H^*$, the equation $Hx=\lambda Gx$ reduces to the
ordinary Hermitian problem for
\[
 C=G^{-1/2}HG^{-1/2},\qquad y=G^{1/2}x.
\]
Consequently every generalized eigenvalue is in $F$, and there is a basis
$x_j$ with $x_j^*Gx_k=\delta_{jk}$. This is useful when the metric itself
contains infinitesimal weights. It does not require a uniform ordinary-real
lower bound on $G$. It does, however, use a square root of an arbitrary
positive definite matrix, and so is one of the statements listed in
\cref{spec:rem:divisibility} as needing a divisible workspace.

Every proof in this section used real closedness of $F$, either through the
splitting of a characteristic polynomial or through a square root of an
arbitrary positive element. \Cref{spec:sec:descent} proves the Hermitian and
normal spectral theorems and the singular-value decomposition again, by a
completely different route, over a Hahn workspace whose value group is
\emph{not} assumed divisible, and with control on which exponents can occur.
Neither proof subsumes the other: the argument here is shorter and applies
to any real closed field, while the argument there applies to any set-sized
ordered abelian value group and returns support information that the
algebraic argument cannot see.

\section{Norms and variational principles with attained extrema}\label{sec:minmax}

\subsection{Operator and Frobenius norms}
Define
\[
 \norm A_{\op}=\sigma_1(A),\qquad
 \norm A_{\Fr}=\sqrt{\tr(A^*A)}
             =\sqrt{\sum_{i,j}|a_{ij}|^2},
\]
with operator norm zero for the zero matrix. The SVD proves, rather than
assumes, that
\begin{equation}\label{eq:opmax}
 \norm A_{\op}=\max_{\norm x=1}\norm{Ax}.
\end{equation}
Indeed, writing $x=Vy$ gives
$\norm{Ax}^2=\sum_j\sigma_j^2|y_j|^2\le\sigma_1^2$, with equality at a
leading right singular vector. Thus the symbol ``max'' is legitimate even
though the sphere was never declared compact.

Both norms are invariant under compatible unitary multiplication. They
satisfy triangle inequalities, and
\begin{gather}\label{eq:normbounds}
 \max_{i,j}|a_{ij}|\le\norm A_{\op}\le\norm A_{\Fr}
     \le\sqrt{mn}\max_{i,j}|a_{ij}|,\\
 \norm{AB}_{\op}\le\norm A_{\op}\norm B_{\op},\qquad
 \norm{AB}_{\Fr}\le\norm A_{\op}\norm B_{\Fr}.
 \label{eq:productbounds}
\end{gather}
For the first inequality in \eqref{eq:normbounds}, apply $A$ to a coordinate
unit vector and then bound one coordinate by the vector norm. The middle
inequality follows from $\sigma_1^2\le\sum_j\sigma_j^2$, and the last
one from the finite sum defining the Frobenius norm. The operator product
bound follows from \eqref{eq:opmax}; the Frobenius bound follows by applying
it to every column of $B$ and summing squares. The companion inequality
$\norm{AB}_{\Fr}\le\norm A_{\Fr}\norm B_{\op}$ follows by taking adjoints.

\subsection{Courant--Fischer and Ky Fan principles}
\begin{theorem}[Variational principles]\label{thm:minmax}
Let $H=H^*$ have eigenvalues $\lambda_1\ge\cdots\ge\lambda_n$. For
$1\le k\le n$,
\begin{align}
 \lambda_k
 &=\max_{\dim W=k}\ \min_{\substack{x\in W\\\norm x=1}}x^*Hx,
 \label{eq:CFmax}\\
 &=\min_{\dim W=n-k+1}\ \max_{\substack{x\in W\\\norm x=1}}x^*Hx.
 \label{eq:CFmin}
\end{align}
Also, for $0\le k\le n$,
\begin{equation}\label{eq:KyFan}
 \sum_{j=1}^k\lambda_j
 =\max_{X^*X=I_k}\tr(X^*HX).
\end{equation}
All inner and outer extrema in these statements are attained.
\end{theorem}
\begin{proof}
For a fixed subspace $W$ with orthonormal basis matrix $X$, the compression
$X^*HX$ is Hermitian. Its largest and smallest eigenvalues are the extrema
of the restricted Rayleigh quotient, by diagonalization. This proves inner
attainment without compactness.

Every $k$-dimensional $W$ meets the span of the eigenvectors for
$\lambda_k,\ldots,\lambda_n$ in a nonzero vector. After normalization that
vector has Rayleigh quotient at most $\lambda_k$. The span of the first
$k$ eigenvectors attains the opposite bound, proving \eqref{eq:CFmax}.
The complementary dimension argument proves \eqref{eq:CFmin}.

For \eqref{eq:KyFan}, express the columns of $X$ in an orthonormal eigenbasis
of $H$. If $a_j$ is the sum of their squared moduli in row $j$, then
$0\le a_j\le1$ and $\sum_j a_j=k$. Hence
\[
 \tr(X^*HX)=\sum_j a_j\lambda_j\le\sum_{j=1}^k\lambda_j.
\]
To see the inequality even for signed $\lambda_j$, subtract $\lambda_k$
from every $\lambda_j$. The terms before $k$ have nonnegative coefficients
and can only decrease on replacing $1$ by $a_j$; the terms after $k$ are
nonpositive. The constant part is unchanged because the weights sum to
$k$. The first $k$ eigenvectors attain equality.
\end{proof}

Applying \eqref{eq:CFmax} to $A^*A$ and taking nonnegative square roots
gives the singular-value principle
\begin{equation}\label{eq:singminmax}
 \sigma_k(A)=\max_{\dim W=k}\ \min_{\substack{x\in W\\\norm x=1}}\norm{Ax}.
\end{equation}
The list is padded with zeros to length $n$ when necessary. Applying
\eqref{eq:KyFan} to $-H$ similarly gives the minimum trace as the sum of the
smallest $k$ eigenvalues. These formulas are algebraic versions of familiar
classical results~\cite{Tao}; their finite-dimensional proofs explain why
they remain valid over $F$.

\subsection{Perturbing eigenvalues and singular values}
\begin{theorem}[Ordered Lipschitz bounds]\label{thm:weyl}
If $H,E$ are Hermitian of the same size, and both eigenvalue lists are
ordered decreasingly, then
\[
 |\lambda_j(H+E)-\lambda_j(H)|\le\norm E_{\op}.
\]
For arbitrary rectangular matrices of the same size,
\begin{equation}\label{eq:singlip}
 |\sigma_j(A+E)-\sigma_j(A)|\le\norm E_{\op}.
\end{equation}
\end{theorem}
\begin{proof}
For a unit vector $x$, Cauchy--Schwarz gives
$|x^*Ex|\le\norm E_{\op}$. Insert the two inequalities for the perturbed
Rayleigh quotient into \eqref{eq:CFmax}, obtaining the Hermitian result.
For singular values, the triangle inequality gives
$\bigl|\norm{(A+E)x}-\norm{Ax}\bigr|\le\norm E_{\op}$ on the unit
sphere. Take the attained minima and maxima in \eqref{eq:singminmax}.
\end{proof}

\begin{theorem}[A Frobenius eigenvalue bound]\label{thm:HW}
For Hermitian $H,G$ with decreasing eigenvalue lists $\lambda,\mu$,
\[
 \sum_{j=1}^n(\lambda_j-\mu_j)^2\le\norm{H-G}_{\Fr}^2.
\]
\end{theorem}
\begin{proof}
Let $u_i,v_j$ be orthonormal eigenbases and set
$p_{ij}=|u_i^*v_j|^2$. The matrix $p$ has nonnegative entries and every row
and column sum is one. Put $c_i=\sum_jp_{ij}\mu_j$. For every $k$,
\[
 \sum_{i=1}^kc_i\le\sum_{j=1}^k\mu_j,
\]
by the weight argument in \cref{thm:minmax}: the column weights
$\sum_{i\le k}p_{ij}$ lie in $[0,1]$ and sum to $k$. Equality holds for
$k=n$. Summation by parts, using $\lambda_i-\lambda_{i+1}\ge0$, now gives
\[
 \tr(HG)=\sum_i\lambda_i c_i\le\sum_i\lambda_i\mu_i.
\]
Expand
$\norm{H-G}_{\Fr}^2=\tr H^2+\tr G^2-2\tr(HG)$ to conclude.
\end{proof}

This proof does not use a decomposition of a doubly stochastic matrix into
permutations, a measure on a unitary group, or a real-valued limiting
argument. It is useful when all the eigenvalue movements, not just the
largest one, matter.

\section{Least squares, pseudoinverses, and best low-rank models}\label{sec:approx}

\subsection{The Moore--Penrose inverse}
Given an SVD $A=U\Sigma V^*$, let $\Sigma^\dagger$ be the transposed
rectangular diagonal matrix whose nonzero entries are $\sigma_j^{-1}$ and
whose remaining entries are zero. Define
\[
 A^\dagger=V\Sigma^\dagger U^*.
\]
Direct diagonal multiplication proves the Penrose identities
\begin{equation}\label{eq:penrose}
 AA^\dagger A=A,\quad A^\dagger AA^\dagger=A^\dagger,
 \quad(AA^\dagger)^*=AA^\dagger,\quad(A^\dagger A)^*=A^\dagger A.
\end{equation}
They characterize $A^\dagger$ uniquely, as in the classical
construction~\cite{Penrose}. For completeness, if $B$ satisfies these
identities, $AB$ is the orthogonal projector onto $\im A$: it is a
Hermitian idempotent with that image. Similarly $BA$ is the orthogonal
projector onto $(\ker A)^\perp$, since $\ker(BA)=\ker A$. The equation
$B=BAB$ then says that $B$ is zero on $(\im A)^\perp$ and, on $\im A$,
is the inverse of the restriction
$A:(\ker A)^\perp\to\im A$. This uniquely determines it and also shows
that the SVD definition is basis-independent.

\begin{theorem}[Least squares with an attained minimum]\label{thm:leastsquares}
For $b\in K^m$, all minimizers of $\norm{Ax-b}$ are
\[
 x=A^\dagger b+z,\qquad z\in\ker A.
\]
The unique minimizer of least norm is $A^\dagger b$, and the minimum residual
is $\norm{(I-AA^\dagger)b}$.
\end{theorem}
\begin{proof}
Write $c=U^*b$ and $y=V^*x$. The squared residual is
\[
 \sum_{j=1}^r|\sigma_j y_j-c_j|^2+\sum_{j=r+1}^m|c_j|^2.
\]
It is minimized precisely by $y_j=c_j/\sigma_j$ for $j\le r$, with
remaining input coordinates free. These free coordinates describe
$\ker A$. They add a nonnegative squared norm to the fixed solution, so
setting them to zero uniquely minimizes the norm. The residual formula is
the corresponding orthogonal projection formula.
\end{proof}

Small singular values are not discarded in this theorem. Their reciprocals
may be infinite, and a mathematically exact least-squares solution may
therefore have infinite norm even when the data are finite. Regularization
in \cref{sec:applications} changes the optimization problem rather than
silently changing what ``nonzero'' means.

\subsection{Low-rank approximation}
Let
\[
 A_k=\sum_{j=1}^k\sigma_j u_jv_j^*,\qquad 0\le k<p.
\]
Here zero terms can be included when $k>r$.

\begin{theorem}[Eckart--Young bounds over $F$]\label{thm:EY}
Among matrices $B$ of rank at most $k$,
\begin{align}
 \min\norm{A-B}_{\op}&=\sigma_{k+1}(A),\label{eq:EYop}\\
 \min\norm{A-B}_{\Fr}^2&=\sum_{j=k+1}^p\sigma_j(A)^2.\label{eq:EYfro}
\end{align}
Both minima are attained by $A_k$.
\end{theorem}
\begin{proof}
The SVD computes the two errors of $A_k$ exactly. For the lower operator
bound, a matrix $B$ of rank at most $k$ has a kernel meeting the span of
$v_1,\ldots,v_{k+1}$ in a nonzero vector. Normalize that vector $x$. Then
\[
 \norm{(A-B)x}=\norm{Ax}\ge\sigma_{k+1}.
\]
For the Frobenius bound, choose $n-k$ orthonormal vectors $w_j$ in $\ker B$
and extend them to an orthonormal basis. Summing squared column norms in
that basis gives
\[
 \norm{A-B}_{\Fr}^2\ge\sum_{j=1}^{n-k}\norm{Aw_j}^2.
\]
The minimum-trace version of \eqref{eq:KyFan}, applied to $A^*A$, bounds
the right side below by the sum of its smallest $n-k$ eigenvalues. That is
exactly \eqref{eq:EYfro}, after padding the singular list with zeros.
\end{proof}

The least-squares lower-rank problem originates in Eckart and
Young~\cite{EY}. Here its optimal error is allowed to be infinitesimal or
infinite. The theorem asserts existence of an optimizer, not uniqueness:
repeated singular values and the operator norm in particular can produce
many optimizers. It also gives a direct operational meaning to every
singular scale: $\sigma_{k+1}$ is the exact error required to replace the
matrix by a rank-$k$ model.

\section{Hahn valuation and singular-value scales}\label{sec:valuation}

From this section onward, unless a result explicitly returns to arbitrary
$F$, work in \eqref{eq:workspace}. For $a\ne0$, put
\[
 v(a)=\min\supp a\in\Gamma,\qquad
 \lc(a)=a_{v(a)}.
\]
Let
\[
 \OO=\{a:v(a)\ge0\},\qquad \mm=\{a:v(a)>0\}.
\]
The residue map $\st:\OO\to\C$ extracts the coefficient at exponent zero.
On real finite elements it is the usual standard part. The valuation of a
matrix or vector means the minimum valuation of its entries; all-zero
objects have valuation $\infty$. We sometimes write $v_{\rm mat}(A)$ to
make this convention explicit.

\begin{remark}[Logical order]\label{spec:rem:logicalorder}
This section and the next use the singular-value decomposition. For a
divisible $\Gamma$ that is \cref{thm:svd}. For an arbitrary set-sized
ordered abelian $\Gamma$ it is \cref{spec:thm:svd}, proved in
\cref{spec:sec:descent} by an argument that uses nothing from this section
onward; likewise the positive square root taken inside \cref{lem:valnorm}
is supplied without real closedness by \cref{spec:lem:normroot}. So the
scale theory below is available over an arbitrary $\Gamma$, and there is no
circularity --- only a forward reference, made so that the classical
development can be read first.
\end{remark}

\subsection{Positivity prevents leading cancellation}
\begin{lemma}[Valuation of a Euclidean norm]\label{lem:valnorm}
For any nonzero finite vector $x$ and any nonzero matrix $A$,
\begin{equation}\label{eq:valnorm}
 v(\norm x)=\min_jv(x_j),\qquad
 v(\norm A_{\op})=v(\norm A_{\Fr})=\min_{i,j}v(a_{ij}).
\end{equation}
\end{lemma}
\begin{proof}
Let $\alpha=\min_jv(x_j)$. The coefficient of $t^{2\alpha}$ in
$\sum_j|x_j|^2$ is
\[
 \sum_{v(x_j)=\alpha}|\lc(x_j)|^2>0.
\]
Thus its valuation is $2\alpha$, and taking its positive square root gives
valuation $\alpha$. Apply the same argument to the entries of $A$ for the
Frobenius norm. Finally \eqref{eq:normbounds} traps the operator norm
between a largest entry modulus and a fixed positive ordinary-real multiple
of it. Both bounds have the same valuation, proving the remaining equality.
\end{proof}

A finite sum of positive leading coefficients is crucial. An arbitrary sum
of complex entries can cancel at its leading exponent; a squared Euclidean
norm cannot.

\begin{lemma}[Unitary matrices preserve the integral lattice]\label{lem:unitaryintegral}
If $U$ is unitary, then $U,U^{-1}\in\OO^{n\times n}$. Its reduction
$\st U$ is an ordinary complex unitary matrix, and
$U\OO^n=\OO^n$.
\end{lemma}
\begin{proof}
Every entry has modulus at most one because its column has squared norm
one. Hence every entry has nonnegative valuation. The inverse is $U^*$ and
has the same property. Reducing $U^*U=I$ gives
$(\st U)^*(\st U)=I$, and the two lattice inclusions follow from integrality
of $U$ and its inverse.
\end{proof}

For the nonzero singular values of $A$, define
\begin{equation}\label{eq:gammas}
 \gamma_j=v(\sigma_j),\qquad
 \gamma_1\le\cdots\le\gamma_r.
\end{equation}
Their order is reversed because a larger positive Hahn number has a smaller
or equal leading exponent. Different singular values may have the same
valuation. The tuple \eqref{eq:gammas} records only scales, not the positive
leading coefficients or the rest of each expansion.

\subsection{The determinantal characterization}
Define the $k$th determinantal valuation by
\begin{equation}\label{eq:Delta}
 \Delta_0(A)=0,\qquad
 \Delta_k(A)=\min_{\substack{I\subseteq\{1,\ldots,m\},\ |I|=k\\
                            J\subseteq\{1,\ldots,n\},\ |J|=k}}
            v(\det A_{I,J}),\quad 1\le k\le p.
\end{equation}
It is $\infty$ exactly when $k>r$.

\begin{lemma}[Integral row and column invariance]\label{lem:DeltaInvariant}
If $L\in\GL_m(\OO)$ and $R\in\GL_n(\OO)$, then
$\Delta_k(LAR)=\Delta_k(A)$ for every $k$.
\end{lemma}
\begin{proof}
By the finite Cauchy--Binet formula, every $k$-minor of $LA$ is a sum of
$k$-minors of $A$ multiplied by $k$-minors of $L$. The latter coefficients
are integral, so every summand has valuation at least $\Delta_k(A)$.
The ultrametric inequality gives $\Delta_k(LA)\ge\Delta_k(A)$. Apply the
same argument to $L^{-1}$ for the reverse inequality. Right multiplication
is identical, or follows by transposition. The argument also covers
vanishing minors.
\end{proof}

\begin{theorem}[Singular scales are determinantal scales]\label{thm:scales}
For a matrix $A$ of rank $r$ and $1\le k\le r$,
\begin{equation}\label{eq:scales}
 \Delta_k(A)=\gamma_1+\cdots+\gamma_k,
 \qquad \gamma_k=\Delta_k(A)-\Delta_{k-1}(A).
\end{equation}
In particular the finite slopes of the determinantal profile are
nondecreasing. The tuple $(\gamma_1,\ldots,\gamma_r)$ is a complete
invariant for equivalence under integral invertible row and column
operations, with the matrix dimensions fixed.
\end{theorem}
\begin{proof}
In an SVD $A=U\Sigma V^*$ both unitary factors and their inverses are
integral by \cref{lem:unitaryintegral}. Thus \cref{lem:DeltaInvariant}
reduces the calculation to $\Sigma$. Its nonzero $k$-minors are products of
$k$ distinct nonzero singular values. Their minimum valuation is the sum of
the $k$ smallest $\gamma_j$, proving \eqref{eq:scales}.

Write $\sigma_j=t^{\gamma_j}u_j$, where $u_j\in\OO^\times$ is a positive
unit. Multiplying diagonal rows by $u_j^{-1}$ shows that $A$ is integrally
equivalent to the rectangular diagonal matrix
$\diag(t^{\gamma_1},\ldots,t^{\gamma_r},0,\ldots)$. Thus two matrices
with the same tuple are equivalent. Conversely, equivalence preserves
$\Delta_k$, hence the tuple and rank.
\end{proof}

For integral matrices this is the singular-value interpretation of their
valuation-ring invariant factors. The valuation ring need not be Noetherian
or a discrete valuation ring. Only finitely many entries and minors are
involved, so a minimum valuation among them exists. Unit equivalence
remembers less than unitary equivalence: the latter preserves all singular
values, not merely their valuations.

An exterior-power interpretation is also useful. Give $\bigwedge^kK^n$
the inner product for which wedges of coordinate basis vectors are
orthonormal. The entries of $\bigwedge^k A$ are the $k$-minors of $A$.
Wedges of the singular bases yield
\[
 \norm{\bigwedge^k A}_{\op}=\sigma_1\cdots\sigma_k.
\]
Applying \cref{lem:valnorm} recovers \eqref{eq:scales}. Thus a determinantal
scale is the scale of the maximal $k$-dimensional volume factor.

\subsection{Why entrywise scales alone are not enough}
\begin{example}[A nearly dependent matrix]\label{ex:near}
Let $q\in\Gamma_{>0}$ and
\[
 A_q=\begin{pmatrix}1&1\\1&1+t^q\end{pmatrix}.
\]
Every entry has valuation zero. Nevertheless,
$\det A_q=t^q$, so $\Delta_1=0$, $\Delta_2=q$, and the singular scales are
$0,q$. In contrast, the matrix of all ones has the same entrywise valuations
and rank one. Changing $q$ changes the hidden singular scale without
changing any of those entrywise valuations.

The matrix $A_q$ is real symmetric and positive definite, since its leading
principal minors are $1$ and $t^q$. Its singular values are its positive
eigenvalues:
\[
 \sigma_\pm=\frac{2+t^q\pm\sqrt{4+t^{2q}}}{2}.
\]
Hence $\lc(\sigma_+)=2$ and
$\sigma_-=(1/2)t^q+\text{terms of valuation greater than }q$.
A minimum-weight assignment using only entry valuations could miss the
cancellation in the determinant. Exact minors or an equivalent elimination
calculation are needed.
\end{example}

\section{Gram coefficients recover scales and leading amplitudes}\label{sec:gram}

The polynomial report already develops Newton profiles of scalar
polynomials~\cite{RepoPoly}. Positive Gram matrices give an especially
clean matrix input to that machinery: their characteristic coefficients
carry no leading cancellation from squared minors.

\begin{theorem}[Positive Cauchy--Binet identity]\label{thm:gram}
For $A\in K^{m\times n}$ write
\begin{equation}\label{eq:gramgen}
 \det(I_n+sA^*A)=\sum_{k=0}^p c_k s^k,
 \qquad c_0=1.
\end{equation}
Then
\begin{equation}\label{eq:gramcoeff}
 c_k=\sum_{\substack{|I|=|J|=k}}|\det A_{I,J}|^2
     =e_k(\sigma_1^2,\ldots,\sigma_r^2).
\end{equation}
For $1\le k\le r$,
\begin{equation}\label{eq:gramval}
 c_k>0,\qquad v(c_k)=2\Delta_k(A).
\end{equation}
For $k>r$, $c_k=0$.
\end{theorem}
\begin{proof}
Expansion of $\det(I+sB)$ by the columns in which an entry of $sB$ is
selected shows that its coefficient of $s^k$ is the sum of the principal
$k$-minors of $B$. For $B=A^*A$, Cauchy--Binet expands the principal minor
indexed by $J$ as
\[
 \det((A^*A)_{J,J})
 =\sum_{|I|=k}\overline{\det A_{I,J}}\det A_{I,J}.
\]
This proves the first equality in \eqref{eq:gramcoeff}. Diagonalizing
$A^*A$ gives the second. If a $k$-minor is nonzero, the leading coefficient
of the sum of its squared moduli and all other squared minors of minimum
valuation is a positive real number. Therefore no leading cancellation is
possible, and the valuation is twice the minimum minor valuation.
\end{proof}

Thus the Gram determinant polynomial provides all the singular scales
without finding an eigenvector. Its nonzero coefficient valuations are
$2\Delta_k$, with successive differences $2\gamma_k$. This statement works
in any ordered group $\Gamma$; it does not require drawing a polygon in an
ordinary real plane.

\subsection{Resolving singular values at a common scale}
Define
\[
 C_k=\st\bigl(t^{-2\Delta_k}c_k\bigr)>0\quad(0\le k\le r),
 \qquad C_0=1.
\]
Let a maximal equal-scale block be
\[
 \gamma_{p_0+1}=\cdots=\gamma_{p_0+d}=\gamma,
\]
where the preceding scales, if any, are strictly smaller and the following
ones, if any, are strictly larger. Write $a_j=\lc(\sigma_j)>0$.

\begin{theorem}[Residual polynomial for a singular-scale block]\label{thm:residual}
For $0\le\ell\le d$,
\begin{equation}\label{eq:residualcoeff}
 \frac{C_{p_0+\ell}}{C_{p_0}}
   =e_\ell(a_{p_0+1}^2,\ldots,a_{p_0+d}^2).
\end{equation}
Consequently the ordinary real polynomial
\begin{equation}\label{eq:residualpoly}
 R_\gamma(X)=X^d+
   \sum_{\ell=1}^d(-1)^\ell\frac{C_{p_0+\ell}}{C_{p_0}}X^{d-\ell}
\end{equation}
has roots precisely the squared leading coefficients in that block,
with multiplicity.
\end{theorem}
\begin{proof}
In the elementary symmetric sum for $c_{p_0+\ell}$, a product of minimum
valuation must contain all the $p_0$ singular values of smaller valuation,
exactly $\ell$ singular values from the block, and none from later blocks.
Indeed, omitting a strictly smaller scale while using a larger one
increases the sum of valuations; the increasing ordering shows that only
choices across the equal boundary can tie the minimum. Therefore its
leading coefficient is
\[
 \left(\prod_{j=1}^{p_0}a_j^2\right)
 e_\ell(a_{p_0+1}^2,\ldots,a_{p_0+d}^2).
\]
For $\ell=0$, the prefactor is exactly $C_{p_0}$. Division proves
\eqref{eq:residualcoeff}, and the elementary symmetric factorization proves
\eqref{eq:residualpoly}.
\end{proof}

This does not give all higher coefficients of the singular values. It gives
their valuations and first nonzero coefficients; repeated roots of
$R_\gamma$ may require another scale refinement. It is already stronger than
knowing only the Newton slopes.

For \cref{ex:near},
\[
 c_1=4+2t^q+t^{2q},\qquad c_2=t^{2q},
\]
so $C_1=4$ and $C_2=1$. The two one-element blocks give leading squares
$4$ and $1/4$, hence leading singular coefficients $2$ and $1/2$.
For a same-scale example,
\[
 A=\diag(2t^\gamma,3t^\gamma),\qquad
 C_1=13,\quad C_2=36,
\]
and $R_\gamma(X)=X^2-13X+36=(X-4)(X-9)$. The valuation alone cannot
distinguish the factors $2$ and $3$, but the positive Gram coefficients do.

\subsection{The Gram case and its parity constraint}
One consequence of \cref{thm:scales} will be used repeatedly below and is
worth recording separately. It is a corollary, not a new identity: it is
what \eqref{eq:gramval} says once the determinantal scales of $A^*A$ are
themselves read off by \eqref{eq:scales}.

\begin{corollary}[Even determinantal scales of Gram matrices]\label{spec:cor:gramscales}
For $A\in K_\Gamma^{m\times n}$ and $1\le k\le\rank(A)$,
\begin{equation}\label{spec:eq:gramscales}
 \Delta_k(A^*A)=2\Delta_k(A).
\end{equation}
\end{corollary}
\begin{proof}
The nonzero eigenvalues of $A^*A$ are the $\sigma_j^2$, of valuations
$2\gamma_j$ in the same nondecreasing order. Apply \eqref{eq:scales} to
$A^*A$ and to $A$.
\end{proof}

Equivalently, the positive Cauchy--Binet identity prevents leading
cancellation in the sums of squared minors, which is why
\eqref{spec:eq:gramscales} has no error term. The parity visible in it is an
arithmetic property of the group $\Gamma$, not a statement about an
integer-valued valuation; \cref{spec:cor:gramcriterion} turns it into an
exact criterion for factoring a positive matrix as a Gram matrix.

\section{Support-certified descent without divisibility}
\label{spec:sec:descent}

Every decomposition in \cref{sec:spectral} was obtained from real closedness
of $F$, which \cref{sec:fields} supplies inside a Hahn workspace only when
$\Gamma$ is divisible. Passing to the divisible hull is harmless for
existence and fatal for arithmetic: it makes every positive square root free,
and so erases the distinction between a matrix whose spectral data need no
new scales and one whose positive root genuinely does. This section keeps the
exact exponent group. Throughout it, $\Gamma$ is an arbitrary set-sized
ordered abelian group and
\[
 F_\Gamma=\R((t^\Gamma)),\qquad K_\Gamma=\C((t^\Gamma))=F_\Gamma(i),
\]
with no divisibility, finite-rank, countability or Archimedean assumption.
Conjugation fixes every $t^\gamma$, and, as before, a positive exponent means
a small monomial.

\subsection{Strong summation is not a limit convention}

\begin{definition}[Strongly summable family]\label{spec:def:summable}
A set-indexed family $(f_i)_{i\in I}$ of Hahn series is \emph{strongly
summable} if $\bigcup_i\supp(f_i)$ is well ordered and, for each exponent
$\gamma$, only finitely many $i$ have a nonzero coefficient at $\gamma$. Its
sum is the coefficientwise finite sum.
\end{definition}

This licenses arbitrary regrouping of a jointly summable family. It does not
declare a sequence of partial sums to converge, and the distinction is not
academic at higher rank.

\begin{example}[A rank-two warning]\label{spec:ex:nonconvergence}
Let $\Gamma=\Z^2$ with the lexicographic order, and put $\delta=(0,1)$ and
$\Omega=(1,0)$, so that $0<n\delta<\Omega$ for every positive integer $n$.
The Hahn identity
\[
 \sum_{n\ge0}t^{n\delta}=(1-t^\delta)^{-1}
\]
is valid, but its partial sums do not converge to the right-hand side in the
valuation topology: the tail after $N$ terms has valuation
$(N+1)\delta<\Omega$ for every $N$. A remainder bound
$v(\text{remainder}_N)\ge N\delta$ is therefore not by itself a convergence
proof in arbitrary rank, and it is never used as one here. This is the same
separation of algebraic identity, strong summability and workspace-relative
topology drawn in \cref{sec:functions}, exhibited in a value group where
even the first two do not bring the third with them.
\end{example}

\subsection{A self-contained support engine}

\Cref{sec:functions} cites the positive-support lemma to Gonshor. Because
finiteness \emph{across all word lengths}, not merely at a fixed degree, is
what the splitting construction needs, we prove it. The combinatorial method
is the minimal-bad-sequence argument associated with Higman's
lemma~\cite{Higman}; the resulting support statement is the Hahn--Neumann
lemma~\cite{Neumann}.

\begin{lemma}[Ordered words]\label{spec:lem:words}
Let $S$ be a well-ordered set, and order finite words in $S$ by subsequence
embedding with coordinatewise increase. Every infinite sequence of words has
a pair $i<j$ such that the $i$th word embeds in the $j$th.
\end{lemma}
\begin{proof}
Suppose a bad sequence exists. Choose it successively so that at each
position the word has the least possible length among words allowing a bad
continuation. No chosen word is empty. Write $w_i=u_ia_i$ with last letter
$a_i$. A sequence in a well order has an infinite nondecreasing subsequence;
choose $a_{i_0}\le a_{i_1}\le\cdots$.

Then $w_0,\ldots,w_{i_0-1},u_{i_0},u_{i_1},\ldots$ is still bad. A
comparison from its initial segment into some $u_{i_j}$ would compare the
corresponding original words; and a comparison $u_{i_j}\preceq u_{i_k}$,
together with $a_{i_j}\le a_{i_k}$, would compare $w_{i_j}$ with $w_{i_k}$.
Both are impossible. This contradicts minimality of the length at position
$i_0$.
\end{proof}

\begin{lemma}[Positive support monoids]\label{spec:lem:neumann}
Let $S\subseteq\Gamma_{>0}$ be well ordered. The additive monoid
$\langle S\rangle_\N$ generated by $S$ is well ordered, and each
$\gamma\in\Gamma$ is the sum of only finitely many finite ordered words with
letters in $S$.
\end{lemma}
\begin{proof}
If the monoid admitted an infinite strictly decreasing sequence, choose a
word representing each term. \Cref{spec:lem:words} gives an embedding of an
earlier word into a later one; positivity of the unused letters and increase
of the matched letters make the earlier sum at most the later one, a
contradiction.

If infinitely many distinct words had sum $\gamma$, apply the same lemma to
those words. Equality of sums forbids any unused positive letter and any
strict increase of a matched letter, so the two words have the same length
and the same letters, contrary to being distinct.
\end{proof}

\begin{lemma}[Hahn evaluation]\label{spec:lem:evaluate}
Let $z_1,\ldots,z_d\in\mm$ and $S=\bigcup_{j=1}^d\supp(z_j)$. Every formal
series $f\in\C[[X_1,\ldots,X_d]]$ has a well-defined strongly summed
evaluation $f(z_1,\ldots,z_d)\in K_\Gamma$, with support in
$M=\langle S\rangle_\N$, and in $M\setminus\{0\}$ when $f(0)=0$. Evaluation
preserves sums, products, conjugation with real formal variables, and formal
composition whose inner series have zero constant coefficient.
\end{lemma}
\begin{proof}
A finite union of well-ordered subsets of a linear order is well ordered.
Every contribution to a coefficient of $f(z)$ determines an ordered word of
exponents in $S$ together with a choice of one of the $d$ variables at each
position. By \cref{spec:lem:neumann} there are, for a fixed exponent, only
finitely many such words and labels, so every coefficient sum is finite and
the support lies in $M$. For a product or a composition, expanding either
side gives the same labelled words with the same finite contributions at
each exponent, so regrouping is legal and the formal identities survive
evaluation. Conjugation preserves supports, and the same coefficientwise
argument applies.
\end{proof}

Matrix formal series reduce to this lemma by taking finitely many scalar
entries as the variables; the entries need not commute, since the scalar
coefficient bookkeeping keeps all ordered matrix products apart. No bound on
the growth of the ordinary complex Taylor coefficients is needed, because a
coefficient at a fixed Hahn exponent is still a finite sum. This is exactly
the mechanism behind \cref{thm:series}, stated once in the generality the
splitting construction requires.

\begin{corollary}[Exponent-preserving inverses and unit roots]\label{spec:cor:unitroots}
If $h\in\mm$, the geometric and binomial series define $(1+h)^{-1}$ and
$(1+h)^a$ for every rational $a$, and all their exponents lie in the group
generated by $\supp(h)$. Every nonzero Hahn series has an inverse without
enlarging its exponent group.
\end{corollary}
\begin{proof}
Apply the formal geometric and binomial identities and
\cref{spec:lem:evaluate}. For a general series factor $f=ct^\alpha(1+h)$;
its inverse is $c^{-1}t^{-\alpha}(1+h)^{-1}$, and both $\alpha$ and the
exponents of $h$ lie in the original support subgroup.
\end{proof}

\begin{lemma}[Formal maps are valuation nonexpanding]\label{spec:lem:lipschitz}
Let $f\in\C[[X_1,\ldots,X_d]]$ and $x,y\in\mm^d$. Then
\[
 v\bigl(f(x)-f(y)\bigr)\ \ge\ \min_jv(x_j-y_j),
\]
and the same entrywise assertion holds for a matrix of formal series.
\end{lemma}
\begin{proof}
There are formal series $g_j\in\C[[X,Y]]$ with
$f(X)-f(Y)=\sum_j(X_j-Y_j)g_j(X,Y)$: for a monomial this is a telescoping
product identity, and collecting by total degree gives the formal identity.
Evaluate with \cref{spec:lem:evaluate}; each $g_j(x,y)$ is integral.
\end{proof}

\begin{lemma}[Norms exist without real closedness]\label{spec:lem:normroot}
For every nonzero $x\in K_\Gamma^n$ we have $x^*x>0$, and $x^*x$ has a
unique positive square root in $F_\Gamma$. The same holds for
$\norm A_{\Fr}^2=\tr(A^*A)$ and for any finite sum of squared moduli of
entries.
\end{lemma}
\begin{proof}
Let $\alpha=v(x)$ and write $x=t^\alpha y$ with $y$ integral and at least
one entry of nonzero residue. The constant coefficient of $y^*y$ is
$\sum_j|\st(y_j)|^2>0$, so $x^*x=ct^{2\alpha}(1+h)$ with $c>0$ real and
$h\in\mm\cap F_\Gamma$. The binomial series $(1+h)^{1/2}$ is a Hahn series
by \cref{spec:lem:evaluate}, so $\sqrt c\,t^\alpha(1+h)^{1/2}\in F_\Gamma$
is the required root, and uniqueness follows from $(u-v)(u+v)=0$ in an
ordered field.
\end{proof}

This is the one place where positivity does all the work. No real closedness
is used, and none is available: if $\gamma\notin2\Gamma$, the positive
element $t^\gamma$ has no square root in $F_\Gamma$ and therefore cannot be
a nonzero sum of squares. That single observation is why an SVD can stay in
the base field while an arbitrary positive matrix square root need not, and
it is quantified exactly in \cref{spec:sec:ramification}. Note also that
\cref{spec:lem:normroot} supplies the root whose existence
\cref{lem:valnorm} takes for granted, so the valuation computation there is
valid over an arbitrary $\Gamma$; \cref{lem:unitaryintegral} likewise used
only integrality and is unaffected.

\subsection{Separated-block lifting}

Let
\begin{equation}\label{spec:eq:D}
 D=\diag(c_1I_{n_1},\ldots,c_sI_{n_s}),\qquad
 c_j\in\R,\quad c_j\ne c_\ell\ (j\ne\ell),
\end{equation}
and write $Q_j$ for the corresponding constant coordinate projections.
Suppose $E=E^*\in\mm^{n\times n}$ and $B=D+E$. Let $S$ be the union of the
supports of the entries of $E$ and put $M=\langle S\rangle_\N$.

\begin{theorem}[Separated-block lifting]\label{spec:thm:split}
There are Hermitian projections $P_1,\ldots,P_s$ and a unitary $W$ over
$K_\Gamma$ such that
\begin{align}
 &P_jP_\ell=0\ (j\ne\ell),\quad \textstyle\sum_jP_j=I,\quad BP_j=P_jB,
   \quad P_j\equiv Q_j\pmod{\mm},\label{spec:eq:projectors}\\
 &W\equiv I\pmod{\mm},\qquad P_jW=WQ_j,\qquad
   W^*BW\text{ is block diagonal relative to }(Q_j).\label{spec:eq:W}
\end{align}
Every entry of $P_j-Q_j$ and of $W-I$ is supported in $M\setminus\{0\}$. The
construction uses universal finite-variable formal series in the entries of
$E$, with ordinary constant coefficients.
\end{theorem}
\begin{proof}
Put $R_0(\zeta)=(\zeta I-D)^{-1}$ and define
\begin{equation}\label{spec:eq:projresolvent}
 P_j(E)=\sum_{k\ge0}\Res_{\zeta=c_j}
    \left[R_0(\zeta)\bigl(ER_0(\zeta)\bigr)^k\right].
\end{equation}
For each $k$ the residue is a homogeneous polynomial of degree $k$ in the
entries of $E$ whose coefficients are ordinary complex numbers, obtained as
residues of rational functions with poles among the fixed real $c_j$. Hence
\cref{spec:lem:evaluate} makes \eqref{spec:eq:projresolvent} a strongly
summed Hahn series with the stated support; the term $k=0$ is $Q_j$.

The projector identities are justified without presupposing a spectral
theorem over $K_\Gamma$. Parametrize Hermitian $E$ by $n^2$ independent
\emph{ordinary real} coordinates. For sufficiently small such coordinates,
disjoint complex circles around the $c_j$ separate the eigenvalues of the
ordinary complex Hermitian matrix $D+E$: if $\lambda$ is farther than
$\norm E$ from every $c_j$, then $I-(\lambda I-D)^{-1}E$ is invertible by
its ordinary geometric series, so $\lambda I-D-E$ is invertible. On those
circles the ordinary resolvent has its Neumann expansion, and the contour
integrals are the orthogonal spectral projectors, as one sees by
diagonalizing the ordinary Hermitian matrix. Their Taylor series at $E=0$
are exactly \eqref{spec:eq:projresolvent}. Therefore $P_j^*=P_j$,
$P_j^2=P_j$, $P_jP_\ell=0$, $\sum_jP_j=I$ and $[B,P_j]=0$ hold as
finite-variable formal identities, every Taylor coefficient of the
difference being zero, and Hahn evaluation transports them to our $E$. This
is a proof of universal coefficient identities by ordinary complex analysis,
\emph{not} contour integration in a surreal topology.

Now set
\begin{equation}\label{spec:eq:polarintertwiner}
 X=\sum_jP_jQ_j,\qquad G=X^*X=\sum_jQ_jP_jQ_j,\qquad W=XG^{-1/2}.
\end{equation}
Here $X\in I+\mm^{n\times n}$ and $G=I+H$ with $H\in\mm^{n\times n}$, and
\begin{equation}\label{spec:eq:Groot}
 G^{-1/2}=\sum_{k\ge0}\binom{-1/2}{k}H^k
\end{equation}
is strongly summable by \cref{spec:lem:evaluate}. It is Hermitian, commutes
with $G$ and squares to $G^{-1}$, so $W^*W=I$; a square matrix with a left
inverse is invertible, so also $WW^*=I$. Orthogonality of the projectors
gives $P_jX=XQ_j$, while $G$, and hence $G^{-1/2}$, commutes with every
$Q_j$; thus $P_jW=WQ_j$. Since $B$ commutes with each $P_j$, its conjugate
$W^*BW$ commutes with every $Q_j$, which is block diagonality. All
nonconstant terms are formal series without constant term in the entries of
$E$, so their supports lie in $M\setminus\{0\}$.
\end{proof}

\begin{remark}[What is canonical and what is not]\label{spec:rem:canonical}
For a fixed residue decomposition $D$ and projections $Q_j$, the formulas
\eqref{spec:eq:projresolvent} and \eqref{spec:eq:polarintertwiner} specify a
particular lift. They do not choose a preferred ordinary eigenbasis inside a
repeated eigenspace of a later leading coefficient; those finite-dimensional
choices are recorded in the splitting tree below. This is the same
distinction as the one after \eqref{eq:spectralprojector}: the eigenspace
projector is canonical, a basis within it is not.
\end{remark}

Writing $E_{j\ell}=Q_jEQ_\ell$ and taking the $k=1$ residues,
\begin{equation}\label{spec:eq:firstorder}
 W=I+K+O_{\mathrm{formal}}(E^2),\qquad
 K_{j\ell}=\frac{E_{j\ell}}{c_\ell-c_j}\ (j\ne\ell),\qquad K_{jj}=0,
\end{equation}
with $K$ skew-Hermitian; $X=I+K+O(E^2)$ gives $X^*X=I+O(E^2)$. The
first-order off-diagonal term of $W^*(D+E)W$ is $E+[D,K]$, which vanishes
between different blocks, and the diagonal block to second degree is
\begin{equation}\label{spec:eq:secondorder}
 c_jI+E_{jj}+\sum_{\ell\ne j}\frac{E_{j\ell}E_{\ell j}}{c_j-c_\ell}
   +O_{\mathrm{formal}}(E^3).
\end{equation}
These $O_{\mathrm{formal}}$ symbols mean support-certified sums of monomials
of the indicated total degrees, not a topological limiting convention. The
second-order coupling correction in \eqref{spec:eq:secondorder} is the
formal-series counterpart of the Schur-complement correction
\eqref{eq:effective}, and \cref{ex:secondorder} is the same phenomenon at
$n=2$.

\begin{proposition}[Local valuation estimate]\label{spec:prop:localstability}
Apply the construction with the same $D,Q_j$ to Hermitian
$E,E'\in\mm^{n\times n}$. Then
\[
 v\bigl(P_j(E)-P_j(E')\bigr)\ge v(E-E'),\qquad
 v\bigl(W(E)-W(E')\bigr)\ge v(E-E').
\]
If $A=\mu I+t^\alpha V(D+E)V^*$ with $V$ constant unitary, and
$v(A'-A)>\alpha$, then the corresponding cluster projectors satisfy
\begin{equation}\label{spec:eq:gapbound}
 v\bigl(\Pi_j(A')-\Pi_j(A)\bigr)\ge v(A'-A)-\alpha.
\end{equation}
\end{proposition}
\begin{proof}
The first pair of inequalities is \cref{spec:lem:lipschitz} applied to the
universal formal entries. For the second, keep the same $\mu,\alpha,V$ and
put $E'=E+t^{-\alpha}V^*(A'-A)V$; its entries are infinitesimal, and
\cref{lem:unitaryintegral} gives $v(E'-E)=v(A'-A)-\alpha$.
\end{proof}

The estimate is sharp already for a two-by-two constant diagonal matrix
perturbed off the diagonal, since \eqref{spec:eq:firstorder} has a nonzero
linear term. Its role is to give a support-compatible \emph{valuation} form
of separated-cluster stability. It does not supersede \cref{thm:subspace}
and \cref{cor:clustergap}, which measure an ordered Frobenius displacement
against an ordered gap; the two are different statements, and
\cref{ex:rotation} shows that an infinitesimal valuation hypothesis alone
cannot control an ordered subspace rotation when the gap is infinitesimal
too.

\subsection{Finite splitting trees and exact spectral descent}

\begin{theorem}[Hermitian Hahn spectral theorem]\label{spec:thm:hermitian}
Let $\Gamma$ be any set-sized ordered abelian group. For every Hermitian
$A\in K_\Gamma^{n\times n}$ there are a unitary $U$ over $K_\Gamma$ and
$\lambda_1,\ldots,\lambda_n\in F_\Gamma$ with
$U^*AU=\diag(\lambda_1,\ldots,\lambda_n)$. For real symmetric input $U$ can
be chosen real orthogonal. Every exponent occurring in $U$ or in the
$\lambda_j$ lies in the subgroup generated by the exponents of the input
entries.
\end{theorem}
\begin{proof}
Induct on $n$, using no algebraic closedness or real closedness assumption.
The case $n=1$ is immediate. Set
\begin{equation}\label{spec:eq:tracesubtract}
 \mu=\tfrac1n\tr(A),\qquad C=A-\mu I.
\end{equation}
If $C=0$ take $U=I$. Otherwise let $\alpha=v(C)$ and write
\begin{equation}\label{spec:eq:normalize}
 t^{-\alpha}C=C_0+E_0,\qquad C_0\in\C^{n\times n},\quad
 E_0\in\mm^{n\times n}.
\end{equation}
The constant matrix $C_0$ is nonzero, Hermitian and traceless, so it has at
least two distinct ordinary real eigenvalues: by the ordinary spectral
theorem a Hermitian matrix with one eigenvalue is scalar, and a traceless
scalar matrix in characteristic zero is zero. Choose an ordinary constant
unitary $V$ with $V^*C_0V=D$ of the form \eqref{spec:eq:D}, with $s\ge2$,
and apply \cref{spec:thm:split} to $D+V^*E_0V$. Then
$W^*V^*AVW=\diag(A_1,\ldots,A_s)$ with each $A_j$ Hermitian of strictly
smaller size. By induction choose unitary $U_j$ diagonalizing $A_j$; then
$U=VW\diag(U_1,\ldots,U_s)$ diagonalizes $A$, and the diagonal entries are
real because the diagonal matrix is Hermitian.

For the support claim let $H$ be the subgroup generated by the input
supports. Both $\mu$ and $C$ have supports in $H$, and $\alpha\in H$, so
normalization only translates supports inside $H$; constant changes of basis
add no exponents; and \cref{spec:thm:split} uses sums of positive exponents
already in $H$, so $W$ and the $A_j$ stay in $K_H$. The induction therefore
runs entirely inside $K_H$. For real input use real constant eigenbases;
every coefficient of the projector and polar constructions is then real.
\end{proof}

This is \cref{thm:spectral} again with two hypotheses removed and one
conclusion added. It does not replace the earlier proof, which is three
lines long and applies to any real closed field; it is the version that
survives when $K_\Gamma$ is not algebraically closed, and it alone says
where the exponents can be.

\begin{corollary}[The splitting tree is finite]\label{spec:cor:tree}
The construction in \cref{spec:thm:hermitian} has at most $n-1$ nontrivial
splitting nodes. This bound is independent of the order type of the input
supports and of the rank or cardinality of $\Gamma$.
\end{corollary}
\begin{proof}
Every internal node has at least two children, and the dimension at a node
is the sum of its children's positive dimensions, so there are at most $n$
leaves. A finite rooted tree in which every internal node has at least two
children has at most one fewer internal node than it has leaves. Termination
follows from the strictly decreasing positive dimensions along each branch.
\end{proof}

At each internal node the construction records the scalar series $\mu$, the
exact leading exponent $\alpha$, the ordinary constant spectral data
$(V,D,Q_j)$, and the well-ordered positive support $S$ of the normalized
perturbation, followed by the inclusion of the lifted projector and unitary
supports in $\langle S\rangle_\N$ and then the certificates of the smaller
blocks. A scalar leaf records its series and multiplicity.

This is a \emph{finite tree of set-valued support data}, not in general a
finite list of monomials: a single node can carry an infinite or uncountable
well-ordered support. Nor is it automatically an effective algorithm from an
arbitrary presentation of a surreal number. Exact zero tests, support access
and extraction of a leading coefficient must be supplied by the
representation --- the same contract as in \cref{sec:algorithm}, now for a
harder task. The coefficient fields remain the full $\R$ and $\C$;
preservation of a smaller coefficient field such as $\Q$ is not asserted.

\begin{example}[A scalar tail beyond every finite truncation]\label{spec:ex:tail}
In the rank-two group of \cref{spec:ex:nonconvergence}, put
\[
 c=\sum_{m\ge1}t^{m\delta},\qquad A=cI+t^\Omega\diag(1,-1).
\]
Every coefficient encountered at the scales $m\delta$ is scalar, so a
procedure that inspects more and more of those coefficients does not reach
the first nonscalar scale after any finite number of inspections. The exact
trace subtraction gives $A-\tfrac12\tr(A)I=t^\Omega\diag(1,-1)$ at once.
This is why \eqref{spec:eq:tracesubtract} is stated as an exact operation,
and it is the spectral counterpart of the truncation failure in
\cref{sec:algorithm}.
\end{example}

\begin{corollary}[Normal spectral descent]\label{spec:cor:normal}
Every normal matrix over $K_\Gamma$ is unitarily diagonalizable over
$K_\Gamma$, without an exponent-group extension.
\end{corollary}
\begin{proof}
Write $A=H+iJ$ with $H=(A+A^*)/2$ and $J=(A-A^*)/(2i)$ Hermitian; normality
is equivalent to $HJ=JH$. Diagonalize $H$ by \cref{spec:thm:hermitian}. The
relation $HJ=JH$ makes the entries of $J$ between distinct exact eigenspaces
of $H$ vanish, so diagonalize the Hermitian restriction of $J$ inside each
block. The resulting unitary diagonalizes $H$ and $J$, hence $A$, and every
operation stayed in the original support group.
\end{proof}

\Cref{prop:normal} proves more about general matrices --- a Schur form for
every square matrix --- but it needs an eigenvalue to exist, hence an
algebraically closed $K$. \Cref{spec:cor:normal} covers only normal matrices
and needs nothing of the sort. Neither contains the other.

\begin{lemma}[A finite-dimensional commutant fact]\label{spec:lem:commutant}
If $A$ is normal over $K_\Gamma$ and $BA=AB$, then $BA^*=A^*B$.
\end{lemma}
\begin{proof}
Use \cref{spec:cor:normal} to make $A$ diagonal. The entry $b_{jk}$ must
vanish when $\lambda_j\ne\lambda_k$, and these are exactly the pairs with
$\overline{\lambda_j}\ne\overline{\lambda_k}$; so $B$ also commutes with the
diagonal conjugate of $A$.
\end{proof}

\begin{theorem}[Finite-witness simultaneous diagonalization]\label{spec:thm:simultaneous}
Let $\mathcal F\subseteq K_\Gamma^{n\times n}$ be a set of pairwise
commuting normal matrices. There is a simultaneous unitary diagonalizer, and
it can be constructed from at most $n-1$ selected members of $\mathcal F$;
consequently its support may be confined to the subgroup generated by those
finitely many members' supports. After selection, one ordinary complex
linear combination of the selected matrices can be chosen whose eigenspaces
are exactly the final joint eigenspaces.
\end{theorem}
\begin{proof}
Begin with the whole space as one block. If every member acts as a scalar
there, stop. Otherwise choose a nonscalar member and decompose the block
into its distinct eigenspaces by \cref{spec:cor:normal}. Every other member
and its adjoint preserve those spaces by \cref{spec:lem:commutant}, so their
restrictions remain normal; repeat on any block on which some member is
nonscalar. Each selection strictly increases the number of joint blocks,
which never exceeds $n$, so at most $n-1$ selections are needed, and the
successive diagonalizers are produced from the selected matrices and
previously constructed bases only.

Let $B_1,\ldots,B_m$ be the selected matrices and let their joint scalar
tuples on the final blocks be $\boldsymbol\lambda^{(a)}\in K_\Gamma^m$;
these are pairwise distinct. For each pair $a\ne b$ the condition
$\sum_jc_j(\lambda_j^{(a)}-\lambda_j^{(b)})=0$ defines a proper
complex-linear subspace of $\C^m$, and a finite union of proper subspaces
over the infinite field $\C$ does not fill $\C^m$. Choose $(c_j)$ outside
that union; then $\sum_jc_jB_j$ has distinct scalar values on the final
blocks, and it is normal because the selected matrices commute with one
another and with one another's adjoints.
\end{proof}

For a family of scalar matrices no member need be selected. This is an
existence theorem for a set-sized family, not a procedure for searching an
arbitrary undecidable one.

\subsection{Singular values stay in the original scale group}

\begin{theorem}[SVD over a nondivisible Hahn field]\label{spec:thm:svd}
Every $A\in K_\Gamma^{m\times n}$ admits $A=U\Sigma V^*$ with $U,V$ unitary
over $K_\Gamma$ and $\Sigma$ rectangular diagonal with entries
$\sigma_1\ge\cdots\ge\sigma_r>0$ in $F_\Gamma$ and zeros elsewhere. The
factors can be chosen with every exponent in the subgroup generated by the
input supports.
\end{theorem}
\begin{proof}
Diagonalize $A^*A$ by \cref{spec:thm:hermitian}. If $v_j$ is a unit
eigenvector with eigenvalue $\lambda_j$, then
$\lambda_j=(Av_j)^*(Av_j)\ge0$. When $\lambda_j>0$,
\cref{spec:lem:normroot} gives $\sigma_j=\sqrt{\lambda_j}\in F_\Gamma$ with
no extension of the group; put $u_j=Av_j/\sigma_j$, which are orthonormal
since $u_j^*u_k=v_j^*A^*Av_k/(\sigma_j\sigma_k)=\delta_{jk}$. When
$\lambda_j=0$, \cref{spec:lem:normroot} forces $Av_j=0$. Extend the nonzero
$u_j$ to an orthonormal basis by finite Gram--Schmidt, selecting standard
basis vectors not yet in the span; every nonzero residual vector has its
norm in $F_\Gamma$ by \cref{spec:lem:normroot}, so no step needs a square
root of an arbitrary positive scalar. Reordering the finitely many positive
singular values gives the stated form and preserves the support field.
\end{proof}

\begin{corollary}[Gram square roots and polar decomposition]\label{spec:cor:gramroot}
For every rectangular $A$ over $K_\Gamma$, the Gram matrix $A^*A$ has its
principal positive semidefinite square root over $K_\Gamma$; in particular
$|A|$ and the polar factorization \eqref{eq:polar} are defined over
$K_\Gamma$ with no enlargement of the value group. For square invertible
$A$, the polar decomposition $A=QP$ has $Q$ unitary and $P$ positive
definite over $K_\Gamma$.
\end{corollary}
\begin{proof}
With the notation of \cref{spec:thm:svd} put
$P=V\diag(\sigma_1,\ldots,\sigma_r,0,\ldots,0)V^*$, so $P^2=A^*A$. In the
invertible square case take $Q=AP^{-1}$, or $Q=UV^*$.
\end{proof}

This does \emph{not} say that every positive semidefinite matrix is a Gram
matrix over a nondivisible base field. The scalar matrix $(t^\gamma)$ with
$\gamma\notin2\Gamma$ is a counterexample, and
\cref{spec:cor:gramcriterion} gives the exact characterization. Note also
that \cref{thm:scales} and \cref{thm:gram} are now available over an
arbitrary $\Gamma$: their proofs used only the SVD, integrality of unitary
factors (\cref{lem:unitaryintegral}), Cauchy--Binet and positivity of a sum
of squared minors, and \cref{spec:thm:svd} supplies the first without
divisibility.

\subsection{The subgroup, and what it is not}

Three sharp examples fix the boundaries of the descent statement.

\begin{example}[New integer combinations, but no new fractions]\label{spec:ex:monoid}
For $\Gamma=\Z$ and $t=t^1$, take
\begin{equation}\label{spec:eq:t3matrix}
 A=\begin{pmatrix}t^3&t^4\\t^4&0\end{pmatrix},
 \qquad
 \lambda_\pm=\frac{t^3}{2}\bigl(1\pm\sqrt{1+4t^2}\bigr),
\end{equation}
so that
$\lambda_+=t^3+t^5-t^7+2t^9-\cdots$ and
$\lambda_-=-t^5+t^7-2t^9+\cdots$. With
$w=(\sqrt{1+4t^2}-1)/(2t)=t-t^3+2t^5-\cdots$, the orthogonal matrix
$U=(1+w^2)^{-1/2}\left(\begin{smallmatrix}1&-w\\w&1\end{smallmatrix}\right)$
satisfies $U^*AU=\diag(\lambda_+,\lambda_-)$, and all exponents are
integers. But the exponent $5$ is \emph{not} in the nonnegative monoid
generated by the raw input exponents $3$ and $4$. A theorem asserting that
eigenvalue supports stay in the raw input monoid would therefore be false.
The correct invariant is the input \emph{subgroup}, supplemented by the
positive monoids of the normalized perturbations at each node.
\end{example}

\begin{example}[An eigenvalue beyond all ordinary perturbation orders]\label{spec:ex:beyond}
In the rank-two group of \cref{spec:ex:nonconvergence} write $t=t^\delta$
and $u=t^\Omega$, and let
$A=\left(\begin{smallmatrix}0&u\\u&t\end{smallmatrix}\right)$. Since
$v(u/t)=\Omega-\delta>0$, the smaller eigenvalue is
\[
 \lambda_-=\frac t2\bigl(1-\sqrt{1+4u^2/t^2}\bigr)
 =-u^2t^{-1}+u^4t^{-3}-2u^6t^{-5}+\cdots,
\]
of valuation $2\Omega-\delta$, which exceeds every finite multiple of
$\delta$. The support stays in $\Z\delta+\Z\Omega$ and the expansion is
certified by the positive generator $2(\Omega-\delta)$. A finite-order
expansion measured only in powers of $t$ never sees this eigenvalue; compare
\cref{ex:three}, where no finite integer cutoff reaches the last singular
direction.
\end{example}

\begin{proposition}[Order-preserving transport]\label{spec:prop:transport}
Let $\phi:\Gamma\hookrightarrow\Lambda$ be an injective order-preserving
group homomorphism. Then
$\phi_*\bigl(\sum_\gamma a_\gamma t^\gamma\bigr)
 =\sum_\gamma a_\gamma s^{\phi(\gamma)}$
is an injective star-preserving field homomorphism, and it preserves
strongly summable families, the support certificates of the splitting
construction, and each chosen spectral or singular-value decomposition.
\end{proposition}
\begin{proof}
An order embedding sends a well-ordered support to a well-ordered support,
and it preserves and reflects equality of sums of exponents, so the finite
coefficient formulas for strong sums and products are unchanged. It
preserves leading coefficients and minima. Hence the trace subtraction, the
normalization, the formal projectors, the polar intertwiner and all selected
ordinary constant eigenbases transport through every node of the finite
tree.
\end{proof}

This lets one symbolic support construction be interpreted at different
surreal scales with the same ordered-group relations. It does not claim that
arbitrary numerical substitutions, in particular ones that identify or
reverse exponents, preserve Hahn summability.

\begin{warning}[Relative degrees can drop in a larger workspace]\label{spec:warn:degreedrop}
Under an embedding into a larger group, an exponent not divisible by $q$ in
the source may become divisible in the target: $\Z\hookrightarrow\Q$ makes
every positive monomial's $q$th root available and kills all ramification.
The spectral identities transport by \cref{spec:prop:transport}, but the
relative ramification group of \cref{spec:sec:ramification} must be
recomputed in the new base group. Degree is preserved under an isomorphism
onto the image group, \emph{not} under arbitrary enlargement of the base
workspace.
\end{warning}

\begin{warning}[This is not a multivariable analytic eigenbasis theorem]\label{spec:warn:crossing}
Consider the real symmetric formal matrix
\begin{equation}\label{spec:eq:crossing}
 H(x,y)=\begin{pmatrix}x&y\\y&-x\end{pmatrix},
\end{equation}
whose characteristic equation is $\lambda^2=x^2+y^2$. There is no solution
$\lambda\in\C[[x,y]]$: its lowest homogeneous term would have to be $ax+by$
with $a^2=b^2=1$ and $2ab=0$. So a general eigenvalue or eigenbasis theorem
in that formal power-series ring is false~\cite{PR,DLLZ}. After the ordered
Hahn specialization $x=t^\alpha$, $y=t^\beta$ with $\beta>\alpha$, by
contrast,
\[
 \sqrt{x^2+y^2}=t^\alpha\sum_{m\ge0}\binom{1/2}{m}t^{2m(\beta-\alpha)}
\]
is a legitimate Hahn series in the original group, because the ordering
permits division by the dominant monomial. Its displayed support may
correspond to negative powers of $x$ in the original two-variable
expression, so it gives neither an element of $\C[[x,y]]$, nor a common
analytic germ through the crossing, nor a global analytic choice of
eigenvectors. The descent theorems resolve the totally ordered scale
problem, and only that one.
\end{warning}

\section{Determinantal ramification and the exact root field}
\label{spec:sec:ramification}

Spectral data need no new scales. Positive matrix roots sometimes do, and
the cost is exactly computable from the invariant already in hand. Every
ordered abelian group is torsion free, so its divisible hull
$\Gamma_\Q=\Gamma\otimes_\Z\Q$ carries a unique order extending that of
$\Gamma$; we write $\gamma/q$ only as an element of this hull unless its
membership in a smaller group has been established. Throughout this section
and the next, $\Gamma$ is again arbitrary.

\subsection{Scalar roots are monomials times base-field units}

\begin{lemma}[The scalar root formula]\label{spec:lem:scalarroot}
Let $a\in F_\Gamma$ be positive, let $\gamma=v(a)$, and let $q\ge2$ be an
integer. There is a unique positive $q$th root in $F_{\Gamma_\Q}$, and it
has the form
\begin{equation}\label{spec:eq:scalarroot}
 a^{1/q}=t^{\gamma/q}u,\qquad u\in F_\Gamma,\quad u>0.
\end{equation}
It belongs to $F_\Gamma$ if and only if $\gamma\in q\Gamma$.
\end{lemma}
\begin{proof}
Factor $a=ct^\gamma(1+h)$ with $c>0$ real and $h\in\mm\cap F_\Gamma$. The
binomial series gives $u=c^{1/q}(1+h)^{1/q}\in F_\Gamma$ by
\cref{spec:cor:unitroots}, with positive constant coefficient, proving
existence and the displayed factorization. The map $x\mapsto x^q$ is
strictly increasing on positive elements of any ordered field, so the
positive root is unique. A root in $F_\Gamma$ must have valuation
$\gamma/q\in\Gamma$, and the formula proves the converse.
\end{proof}

\begin{lemma}[Finite-index Hahn extensions]\label{spec:lem:finiteindex}
Let $\Gamma\subseteq\Lambda\subseteq\Gamma_\Q$ with $\Lambda/\Gamma$ finite.
Then
\begin{equation}\label{spec:eq:finiteindex}
 [K_\Lambda:K_\Gamma]=[F_\Lambda:F_\Gamma]=|\Lambda/\Gamma|,
\end{equation}
and a choice of coset representatives $\eta_h$ gives a basis
$(t^{\eta_h})_{h\in\Lambda/\Gamma}$ over the respective base field. If
$q\Lambda\subseteq\Gamma$, the complex extension is Galois and
\begin{equation}\label{spec:eq:galois}
 \Gal(K_\Lambda/K_\Gamma)=\Hom(\Lambda/\Gamma,\mu_q),
\end{equation}
where $\mu_q\subset\C^\times$ is the group of $q$th roots of unity and a
character $\chi$ acts by $t^\eta\mapsto\chi(\eta+\Gamma)t^\eta$. The
extension is totally ramified: its residue field is still $\C$ and its
ramification index is the full degree.
\end{lemma}
\begin{proof}
Split the support of a Hahn series into the finitely many cosets of
$\Gamma$. Translating the part in coset $h$ by $-\eta_h$ gives a
well-ordered subset of $\Gamma$, so every series has a unique expression
$f=\sum_h t^{\eta_h}f_h$ with $f_h\in K_\Gamma$, and the analogous real
expression; uniqueness holds because supports in distinct cosets are
disjoint. This proves the basis and the degree.

Multiplication by a character on each coset preserves all coefficientwise
sums and products, so it defines an automorphism. A finite abelian group
killed by $q$ has exactly its own order many characters into $\mu_q$, and
these separate points, as one sees by decomposing it into cyclic factors.
We have thus exhibited as many distinct automorphisms as the degree, so the
extension is Galois with the stated group. The residue field is unchanged
and the value group grows by the full index, which is total ramification.
\end{proof}

The finite-index hypothesis is essential, not cosmetic. For an infinite
group extension, the field generated by all the new monomials need not
contain every Hahn series supported in the enlarged group. We invoke
\cref{spec:lem:finiteindex} only for extensions generated by finitely many
fractional exponents of bounded denominator.

\subsection{Existence, uniqueness, and the exact extension}

\begin{theorem}[Principal-root field]\label{spec:thm:rootfield}
Let $A\in K_\Gamma^{n\times n}$ be Hermitian positive semidefinite, with
positive eigenvalues $\lambda_1,\ldots,\lambda_r$ counted with multiplicity,
and put $\gamma_j=v(\lambda_j)$ and
$\Lambda=\Gamma+\sum_{j=1}^r\Z\gamma_j/q$. There is a unique positive
semidefinite Hermitian $B$ over $K_{\Gamma_\Q}$ with $B^q=A$. It lies in
$K_\Lambda^{n\times n}$, and
\begin{equation}\label{spec:eq:exactrootfield}
 K_\Gamma(B)=K_\Lambda,
\end{equation}
where $K_\Gamma(B)$ denotes the field obtained by adjoining all entries
of $B$.
\end{theorem}
\begin{proof}
Choose a unitary diagonalization over the base field, furnished by
\cref{spec:thm:hermitian}:
$A=U\diag(\lambda_1,\ldots,\lambda_r,0,\ldots,0)U^*$. The diagonal entries
are nonnegative, being values of the quadratic form on unit eigenvectors. By
\cref{spec:lem:scalarroot}, $b_j=\lambda_j^{1/q}=t^{\gamma_j/q}u_j$ with
$u_j\in F_\Gamma$ positive. Set
\begin{equation}\label{spec:eq:principalroot}
 B=U\diag(b_1,\ldots,b_r,0,\ldots,0)U^*,
\end{equation}
which is positive semidefinite, satisfies $B^q=A$, and has entries in
$K_\Lambda$.

For uniqueness let $C=C^*\succeq0$ satisfy $C^q=A$ in the enlarged field.
Then $C$ and $A$ commute, so \cref{spec:thm:simultaneous} gives a
simultaneous Hermitian diagonalization there. On an eigenspace of $A$ with
value $\lambda$, every eigenvalue of $C$ is nonnegative with $q$th power
$\lambda$, and the ordered scalar uniqueness of
\cref{spec:lem:scalarroot} forces it to be the chosen root, including $0$
when $\lambda=0$. Hence $C=B$.

Let $L=K_\Gamma(B)$. Since $U$ lies in the base field, the entries of
$U^*BU$ show every $b_j\in L$; dividing by the base-field unit $u_j$ gives
$t^{\gamma_j/q}\in L$. These finitely many monomials, together with
monomials already in the base field, generate all coset representatives of
$\Lambda/\Gamma$ by multiplication and inversion, and the quotient is
finite, being generated by $r$ elements killed by $q$. The basis in
\cref{spec:lem:finiteindex} therefore gives $K_\Lambda\subseteq L$, and the
reverse inclusion is \eqref{spec:eq:principalroot}.
\end{proof}

This is the exact sense in which \cref{prop:positive} needs a divisible
workspace. The positive semidefinite root exists and is unique, but over a
general $\Gamma$ it lives in $K_\Lambda$, and $\Lambda$ can be strictly
larger than $\Gamma$.

\subsection{Replacing eigenvalues by minors}

The eigenvalues in \cref{spec:thm:rootfield} can be eliminated entirely. By
\eqref{eq:scales} the determinantal valuations and the eigenvalue valuations
determine one another, so the field of definition is computable from
finitely many minors and no spectral decomposition at all.

\begin{theorem}[Determinantal ramification law]\label{spec:thm:ramification}
Let $A\in K_\Gamma^{n\times n}$ be Hermitian positive semidefinite of rank
$r$, let $\Delta_k=\Delta_k(A)$ be as in \eqref{eq:Delta}, and for an
integer $q\ge2$ put
\begin{equation}\label{spec:eq:rootgroup}
 \Gamma_{A,q}=\Gamma+\sum_{k=1}^r\Z\frac{\Delta_k}q\subseteq\Gamma_\Q,
 \qquad
 R_q(A)=\bigl\langle \Delta_1+q\Gamma,\ldots,\Delta_r+q\Gamma\bigr\rangle
   \subseteq\Gamma/q\Gamma.
\end{equation}
Then the principal positive semidefinite root satisfies
\begin{align}
 K_\Gamma(A^{1/q})&=K_{\Gamma_{A,q}},\label{spec:eq:fieldmain}\\
 [K_\Gamma(A^{1/q}):K_\Gamma]&=|R_q(A)|\le q^r.\label{spec:eq:degreemain}
\end{align}
In particular
\begin{equation}\label{spec:eq:criterionmain}
 A^{1/q}\in K_\Gamma^{n\times n}
 \quad\Longleftrightarrow\quad
 \Delta_k(A)\in q\Gamma\quad(1\le k\le r).
\end{equation}
The extension is finite, abelian and totally ramified, with Galois group
isomorphic to $\Hom(R_q(A),\mu_q)$.
\end{theorem}
\begin{proof}
Order the $\gamma_j$ nondecreasingly. By \eqref{eq:scales},
$\Delta_k=\gamma_1+\cdots+\gamma_k$, so
\begin{equation}\label{spec:eq:samelattice}
 \Gamma+\sum_j\Z\frac{\gamma_j}q=\Gamma+\sum_k\Z\frac{\Delta_k}q:
\end{equation}
each $\Delta_k/q$ is a sum of the $\gamma_j/q$, and each
$\gamma_j/q=(\Delta_j-\Delta_{j-1})/q$. The map
$\Gamma/q\Gamma\to(q^{-1}\Gamma)/\Gamma$ sending $\gamma+q\Gamma$ to
$\gamma/q+\Gamma$ is an isomorphism, and it identifies $R_q(A)$ with
$\Gamma_{A,q}/\Gamma$. Combining \cref{spec:thm:rootfield} with
\cref{spec:lem:finiteindex} gives
\eqref{spec:eq:fieldmain}, \eqref{spec:eq:degreemain} and the Galois
assertion, and the extension is trivial exactly when every generating class
vanishes, which is \eqref{spec:eq:criterionmain}.
\end{proof}

Two things deserve emphasis. First, \eqref{spec:eq:degreemain} is not merely
a bound: it computes the exact degree from finitely many minor valuations,
without computing an eigenbasis or the eigenvalues. Second, the same
invariant $\Delta_k$ that \cref{thm:scales} uses to record singular scales
here controls a finite radical extension; the determinantal profile is doing
two different jobs, and \eqref{spec:eq:criterionmain} is the second one in
its sharpest form.

\begin{corollary}[Exact Gram-factorization criterion]\label{spec:cor:gramcriterion}
A Hermitian positive semidefinite $A$ over $K_\Gamma$ can be written
$A=C^*C$ with $C$ over $K_\Gamma$ if and only if
$\Delta_k(A)\in2\Gamma$ for $1\le k\le\rank(A)$. The condition permits $C$
to be square, and it is necessary for every rectangular Gram
representation.
\end{corollary}
\begin{proof}
Necessity is \cref{spec:cor:gramscales}: if $A=C^*C$ then
$\Delta_k(A)=2\Delta_k(C)\in2\Gamma$. For sufficiency,
\cref{spec:thm:ramification} with $q=2$ supplies the principal square root
$B$ over $K_\Gamma$; take $C=B$.
\end{proof}

\begin{corollary}[Rational powers]\label{spec:cor:rationalpower}
Let $A$ be positive definite, $p\in\Z$, $q\ge2$, $\gcd(p,q)=1$. Then the
principal power $A^{p/q}$ generates the same extension as $A^{1/q}$. For
positive semidefinite $A$ the same holds when $p>0$.
\end{corollary}
\begin{proof}
The nonzero diagonal spectral entries of $A^{p/q}$ are $b_j^p$. If
$ap+bq=1$, then $b_j=(b_j^p)^a\lambda_j^b$, so adjoining these entries is
field-theoretically equivalent to adjoining the $b_j$. Zero eigenvalues are
omitted, and negative powers require invertibility.
\end{proof}

\subsection{Prime and composite denominators}

When $q$ is prime, $\Gamma/q\Gamma$ is a vector space over $\mathbf F_q$ and
\[
 [K_\Gamma(A^{1/q}):K_\Gamma]
 =q^{\dim_{\mathbf F_q}\Span\{\Delta_k+q\Gamma:1\le k\le r\}}.
\]
For composite $q$ the exact invariant is the finite subgroup, not the rank
of a vector space: a single class of order $2$ modulo $4\Gamma$ contributes
a quadratic extension, not a quartic one.

If $\Gamma=\Z^d$ with any fixed compatible order and the $\Delta_k$ are
integer vectors, put $L=q\Z^d+\sum_{k=1}^r\Z\Delta_k\subseteq\Z^d$. Then
\begin{equation}\label{spec:eq:latticedegree}
 |R_q(A)|=\frac{q^d}{[\Z^d:L]},
\end{equation}
and the index is computed by the Smith normal form of the integer matrix
with columns $qI_d,\Delta_1,\ldots,\Delta_r$. The root-field degree is
therefore a finite integer-lattice computation once the minor valuations are
known exactly, and no discreteness in the \emph{order} of $\Gamma$ is needed
for it.

\subsection{Three examples at the boundary}

\begin{example}[A positive matrix whose square root genuinely ramifies]\label{spec:ex:ramifying}
Take
\begin{equation}\label{spec:eq:positive2}
 A=\begin{pmatrix}1+t&1\\1&1\end{pmatrix}\in\C((t))^{2\times2},
\end{equation}
positive definite because $x^*Ax=|x_1+x_2|^2+t|x_1|^2>0$ for $x\ne0$. The
minor minima are $\Delta_1=0$ and $\Delta_2=v(t)=1$, so the principal square
root generates exactly $\C((t^{\frac12\Z}))$, a quadratic extension of
$\C((t^\Z))$. Explicitly
\begin{equation}\label{spec:eq:positive2root}
 B=\frac{A+\sqrt t\,I}{\sqrt{2+t+2\sqrt t}},
\end{equation}
whose denominator has a positive constant leading term and so introduces no
exponent beyond $\tfrac12\Z$; Cayley--Hamilton gives $B^2=A$, and positivity
of numerator and denominator gives $B\succeq0$. Moreover
$T=\tr(B)=\sqrt{2+t+2\sqrt t}$ and $\sqrt t=(T^2-2-t)/2$, which makes the
primitive-trace phenomenon of \cref{spec:sec:trace} directly visible.
\end{example}

\begin{example}[A Gram witness can use a larger group than its product]\label{spec:ex:gramwitness}
With $t=t^\delta$, $u=t^\Omega$ in the rank-two group of
\cref{spec:ex:nonconvergence}, put
\[
 C=\begin{pmatrix}t&1\\0&u\end{pmatrix},\qquad
 A=C^*C=\begin{pmatrix}t^2&t\\t&1+u^2\end{pmatrix}.
\]
Relative to $\Gamma=\Z\delta+\Z\Omega$ the minor valuations are
$\Delta_1=0$ and $\Delta_2=2\delta+2\Omega$, so the square root needs no
extension while the fourth root generates one of degree $2$. But the exact
input support group of $A$ \emph{alone} is $H=\Z\delta+2\Z\Omega$, and
relative to $H$ the value $2\delta+2\Omega$ is not in $2H$, so $A^{1/2}$
generates a quadratic extension of $K_H$ and recovers the missing scale $u$.
This does not contradict \cref{spec:cor:gramroot}: the displayed Gram
witness $C$ already uses that scale. It is why every statement here names
both the input matrix and the base group, and why
\cref{spec:prop:exactlocalization} is worth distinguishing from
\cref{prop:localization}.
\end{example}

\begin{example}[The degree bound is attained in every rank]\label{spec:ex:sharpdegree}
Let $\Gamma=\Z^r$ with the lexicographic order and standard basis
$e_1,\ldots,e_r$, each positive, and let
$A=\diag(t^{e_1},\ldots,t^{e_r})$. The nonzero eigenvalue valuations
generate all of $\Gamma/q\Gamma$, and their partial sums, which are the
determinantal values after ordering, generate the same subgroup. Hence
$[K_\Gamma(A^{1/q}):K_\Gamma]=q^r$. This rules out any degree bound
depending only on $q$ and not on the rank.
\end{example}

\section{The trace is a primitive generator}
\label{spec:sec:trace}

The extension of \cref{spec:thm:ramification} is generated by the $n^2$
entries of $A^{1/q}$. It is generated by one of their linear functionals.

\subsection{A positive-coset lemma}

\begin{lemma}[Positive radical sums retain every generator]\label{spec:lem:positivetrace}
Let $a_1,\ldots,a_r\in F_\Gamma$ be positive, let $q\ge2$, and write
$b_j=a_j^{1/q}=t^{\gamma_j/q}u_j$ as in \eqref{spec:eq:scalarroot}. Put
$\Lambda=\Gamma+\sum_j\Z\gamma_j/q$. Then for any strictly positive
$w_j\in F_\Gamma$,
\begin{equation}\label{spec:eq:positivesum}
 K_\Gamma\Bigl(\sum_{j=1}^rw_jb_j\Bigr)=K_\Lambda.
\end{equation}
\end{lemma}
\begin{proof}
Let $G=\Lambda/\Gamma$ and $h_j=\gamma_j/q+\Gamma$; these classes generate
$G$. Choose a representative $\eta_h$ for each class occurring among the
$h_j$. The sum $T=\sum_jw_jb_j$ has the coset decomposition
\[
 T=\sum_ht^{\eta_h}s_h,\qquad
 s_h=\sum_{j:h_j=h}w_j\,t^{\gamma_j/q-\eta_h}u_j\in F_\Gamma,
\]
and every summand in $s_h$ is positive in the ordered field $F_\Gamma$, so
$s_h>0$ and in particular $s_h\ne0$ for each occurring class.

By \cref{spec:lem:finiteindex} every automorphism of $K_\Lambda/K_\Gamma$ is
a character $\chi:G\to\mu_q$. If such an automorphism fixes $T$, then
\[
 0=\chi(T)-T=\sum_h(\chi(h)-1)t^{\eta_h}s_h.
\]
Distinct cosets have disjoint supports, so every coefficient $\chi(h)-1$
vanishes; since the occurring $h_j$ generate $G$, the character is trivial.
Thus $T$ has trivial stabilizer in the Galois group, and the fundamental
correspondence gives $[K_\Gamma(T):K_\Gamma]=|G|$, which is the asserted
equality.
\end{proof}

\begin{theorem}[Primitive trace]\label{spec:thm:primitivetrace}
Under the hypotheses of \cref{spec:thm:ramification},
\begin{equation}\label{spec:eq:tracemain}
 K_\Gamma\bigl(\tr(A^{1/q})\bigr)=K_\Gamma(A^{1/q})=K_{\Gamma_{A,q}}.
\end{equation}
More generally, if $\lambda_1,\ldots,\lambda_r$ are the positive eigenvalues
of $A$, then every sum $\sum_ja_j\lambda_j^{1/q}$ with $a_j\in F_\Gamma$ and
$a_j>0$ is a primitive generator of the same field.
\end{theorem}
\begin{proof}
For the principal root \eqref{spec:eq:principalroot}, trace invariance gives
$\tr(B)=\sum_{j=1}^r\lambda_j^{1/q}$. Apply
\cref{spec:lem:positivetrace} with $w_j=1$, and use
\cref{spec:thm:rootfield} together with \eqref{spec:eq:samelattice}. The
weighted statement is the same lemma with arbitrary strictly positive
$w_j$.
\end{proof}

The proof identifies a mechanism rather than invoking an abstract
primitive-element theorem. In a general finite extension most linear
combinations may be primitive while special choices fail; here the
\emph{specific} unweighted trace always succeeds, because positivity makes
its coset support a generating set. Strictly positive coefficients are the
whole content of the hypothesis, as \cref{spec:warn:positivity} shows.

\begin{corollary}[The corresponding real field]\label{spec:cor:realtrace}
With $T=\tr(A^{1/q})$ and $\Lambda=\Gamma_{A,q}$, we have
$F_\Gamma(T)=F_\Lambda$.
\end{corollary}
\begin{proof}
Both fields are ordered and $T\in F_\Lambda$. Adjoining $i$ has degree $2$
over each of these real fields, so the tower formula applied to
$K_\Gamma(T)=K_\Lambda$ gives
$[F_\Gamma(T):F_\Gamma]=[K_\Lambda:K_\Gamma]=|\Lambda/\Gamma|$, which is
$[F_\Lambda:F_\Gamma]$ by \cref{spec:lem:finiteindex}.
\end{proof}

\subsection{An explicit irreducible orbit polynomial}

Let $b_j=\lambda_j^{1/q}$, $G=\Gamma_{A,q}/\Gamma$ and $h_j=v(b_j)+\Gamma$.
The minimal polynomial of the trace over $K_\Gamma$ is
\begin{equation}\label{spec:eq:orbitpoly}
 \mathcal P_{A,q}(X)=\prod_{\chi\in\Hom(G,\mu_q)}
   \Bigl(X-\sum_{j=1}^r\chi(h_j)b_j\Bigr).
\end{equation}
Its coefficients lie in $K_\Gamma$ because the Galois group permutes the
factors, and the factors are distinct by \cref{spec:lem:positivetrace}, so
the polynomial is irreducible of degree $|R_q(A)|$. It is a formula inside a
finite extension, not necessarily a fast coefficient-extraction algorithm.

\begin{example}[A biquadratic field encoded in one positive trace]\label{spec:ex:biquadratic}
Let $\Gamma=\Z\delta\oplus\Z\Omega$ with $0<n\delta<\Omega$ for every
positive integer $n$, and let $A=\diag(t^\delta,t^\Omega)$ with $q=2$. The
minor valuations are $\Delta_1=\delta$ and $\Delta_2=\delta+\Omega$, and
their classes span a subgroup of order $4$ in $\Gamma/2\Gamma$. So
$B=A^{1/2}$ generates a biquadratic extension, and so does
$T=t^{\delta/2}+t^{\Omega/2}$. Writing $t=t^\delta$ and $u=t^\Omega$, the
minimal polynomial is
\begin{equation}\label{spec:eq:biquadratic}
 X^4-2(t+u)X^2+(t-u)^2,
\end{equation}
whose degree $4$ is certified by the two independent exponent classes and
not by an unproved irreducibility guess about the displayed polynomial.
\end{example}

\begin{warning}[Positivity cannot be dropped]\label{spec:warn:positivity}
For $A=t^\delta I_2$ with $\delta\notin2\Gamma$, the matrix
$C=\diag(t^{\delta/2},-t^{\delta/2})$ is a Hermitian square root generating
a quadratic extension, and $\tr(C)=0$. It is not the positive root.
\Cref{spec:thm:primitivetrace} concerns the principal positive root and
strictly positive weights, never an arbitrary branch or a signed sum.
\end{warning}

\subsection{Several positive roots at once}

\begin{corollary}[Joint ramification and a joint primitive element]\label{spec:cor:joint}
Let $A_1,\ldots,A_m$ be positive semidefinite Hermitian matrices over
$K_\Gamma$, of possibly different sizes, and let $q_\ell\ge2$. Put
\[
 \Lambda=\Gamma+\sum_{\ell=1}^m\sum_{k=1}^{\rank(A_\ell)}
   \Z\frac{\Delta_k(A_\ell)}{q_\ell}.
\]
Then the field generated by all entries of all $A_\ell^{1/q_\ell}$ is
$K_\Lambda$, of degree $|\Lambda/\Gamma|$, and the single scalar
$T=\sum_\ell\tr(A_\ell^{1/q_\ell})$ is a primitive generator. No commutation
assumption on the $A_\ell$ is needed.
\end{corollary}
\begin{proof}
The compositum of the individual monomial extensions contains the monomials
generating $\Lambda$, and the quotient is finite and killed by
$Q=\lcm(q_1,\ldots,q_m)$, so \cref{spec:lem:finiteindex} identifies the
compositum with $K_\Lambda$. Decompose the finitely many positive scalar
spectral roots occurring in the traces into cosets of $\Gamma$; every
occurring coefficient sum is positive and the occurring classes generate
$\Lambda/\Gamma$, so the character argument of
\cref{spec:lem:positivetrace}, now with $\mu_Q$, proves primitivity.
\end{proof}

\subsection{Back inside the surcomplex numbers}

\begin{corollary}[Exact-support surcomplex spectral calculus]\label{spec:cor:surcomplex}
Let $A$ be a finite surcomplex matrix and let $H$ be its exact input support
group, as in \cref{spec:prop:exactlocalization}. If $A$ is Hermitian, it has
real surreal eigenvalues and a unitary eigenbasis whose normal-form
exponents stay in $H$; the same support assertion holds for normal
diagonalization and for the SVD of an arbitrary finite surcomplex matrix.
If $A$ is in addition positive semidefinite, then for every $q\ge2$
\[
 K_H\bigl(A^{1/q}\bigr)
   =K_{H+\sum_k\Z\Delta_k(A)/q}
   =K_H\bigl(\tr(A^{1/q})\bigr),
\]
and the degree over $K_H$ is the order of the subgroup generated by the
classes $\Delta_k(A)+qH$ in $H/qH$.
\end{corollary}
\begin{proof}
Apply \cref{spec:thm:hermitian,spec:cor:normal,spec:thm:svd,%
spec:thm:ramification,spec:thm:primitivetrace} inside the set-sized field
$K_H$, then transport by the normal-form embedding of
\cref{spec:prop:exactlocalization}. The positive branch agrees with the
ambient surreal positive root by ordered-field uniqueness together with the
uniqueness proof in \cref{spec:thm:rootfield}.
\end{proof}

The ambient surreal exponent group is divisible, so the full class field has
room for all of these fractional scales; the content of the corollary is not
that positive roots fail to exist in $\SC$. It is their exact \emph{relative
cost} over the smallest Hahn workspace containing the data --- a cost that
becomes invisible the moment one replaces $H$ by its divisible hull, which
is precisely what \cref{prop:localization} does.

\section{Rank at a chosen scale}\label{sec:rank}

\subsection{Exact rank and standard-part rank}
\begin{corollary}[Rank after reduction]\label{cor:stRank}
If $A\in\OO^{m\times n}$ has nonzero singular scales
$\gamma_1,\ldots,\gamma_r$, then every $\gamma_j\ge0$ and
\[
 \rank_\C(\st A)=\#\{j:\gamma_j=0\}.
\]
\end{corollary}
\begin{proof}
By \cref{lem:valnorm}, $v(\sigma_1)=v(A)\ge0$, so all singular values
are integral. Reduce an SVD. Its unitary factors reduce to invertible
ordinary unitary matrices, and its diagonal entries reduce to nonzero
numbers exactly at valuation zero.
\end{proof}

An invertible matrix can therefore have a singular standard part. Conversely,
an invertible standard part certifies invertibility of an integral square
matrix. The inverse need not be integral in the first situation, but is
integral in the second: the latter is exactly the case in which all its
singular scales are zero.

\subsection{An intrinsic input filtration}
Fix $\theta\in\Gamma$ and define an $\OO$-submodule of the input lattice by
\begin{equation}\label{eq:latticefiltration}
 M_\theta(A)=\{x\in\OO^n:Ax\in t^\theta\OO^m\}.
\end{equation}
Its image $\st(M_\theta(A))$ under coordinatewise reduction is a
$\C$-linear subspace of $\C^n$. Also define
\[
 M_{>\theta}(A)=\{x\in\OO^n:v(Ax)>\theta\}.
\]
These definitions work even when $A$ or $t^\theta$ is nonintegral; only the
input vectors are required to lie in the integral lattice.

\begin{theorem}[Filtration dimensions]\label{thm:filtration}
For every $\theta\in\Gamma$,
\begin{align}
 \dim_\C\st(M_\theta(A))
   &=n-\#\{j:\gamma_j<\theta\},\label{eq:filtrationweak}\\
 \dim_\C\st(M_{>\theta}(A))
   &=n-\#\{j:\gamma_j\le\theta\}.\label{eq:filtrationstrict}
\end{align}
Only the nonzero singular values enter the counts.
\end{theorem}
\begin{proof}
Use an SVD $A=U\Sigma V^*$. Both unitary factors preserve integral
lattices and their reductions are invertible. Write $x=Vy$. For $j\le r$,
condition \eqref{eq:latticefiltration} becomes
\[
 v(y_j)\ge\max(0,\theta-\gamma_j).
\]
The kernel coordinates $j>r$ are unrestricted in $\OO$. The reduction of
$y_j$ must be zero exactly when $\gamma_j<\theta$; all other reduced
coordinates can be chosen arbitrarily. This gives
\eqref{eq:filtrationweak}. For the strict condition, reduction of $y_j$
vanishes exactly when $\gamma_j\le\theta$. Multiplication by the
invertible $\st V$ preserves the dimension, proving
\eqref{eq:filtrationstrict}.
\end{proof}

The two formulas deliberately differ at the boundary. A direction of
singular scale exactly $\theta$ remains visible when the output threshold
is strict, and disappears if only outputs of strictly larger magnitude are
counted. It is helpful to fix explicit notation:
\[
 r_{<\theta}(A)=\#\{j:\gamma_j<\theta\},\qquad
 r_{\le\theta}(A)=\#\{j:\gamma_j\le\theta\}.
\]
These are not replacements for algebraic rank. They describe how much rank
is visible at a specified valuation resolution. In the integral case,
$r_{\le0}(A)=\rank(\st A)$.

\begin{example}[Three genuinely different levels]\label{ex:three}
Choose a workspace containing $1$ and $\omega$ in its value group and let
\[
 A=\diag(1,t^2,t^\omega),\qquad t=\omega^{-1}.
\]
Its exact rank is three and its standard-part rank is one. Its closed-scale
rank $r_{\le\theta}$ is one for $0\le\theta<2$, two for
$2\le\theta<\omega$, and three for $\theta\ge\omega$.
No finite integer cutoff $N$ reaches the last direction. This is not a
statement about a slow numerical convergence rate; it is an ordered-group
fact, $N<\omega$ for every ordinary $N$.
\end{example}

\section{A root-free scale algorithm and its computational contract}\label{sec:algorithm}

\subsection{Minimum-valuation elimination}
The determinantal formula is finite, but enumerating all minors is
expensive. A valuation-aware elimination computes the same scale tuple
without singular vectors or algebraic root extraction.

\begin{proposition}[Exact scale elimination]\label{prop:elimination}
Suppose exact field arithmetic, exact zero testing, and comparison of
valuations of nonzero elements are available. The following algorithm
terminates after at most $\min(m,n)$ pivots and returns
$\gamma_1,\ldots,\gamma_r$ in nondecreasing order.

In the active submatrix, choose a nonzero entry of minimum valuation and
move it to the leading position by row and column swaps. Clear the remaining
entries in its column by row operations, and then clear its row by column
operations. Record its valuation and recurse on the trailing submatrix.
Stop when the active submatrix is zero.
\end{proposition}
\begin{proof}
Write one active block as
\[
 \begin{pmatrix}a&b\\c&D\end{pmatrix},\qquad
 v(a)=\alpha\le v(\text{every entry}).
\]
The row multipliers $c_i/a$ and column multipliers $b_j/a$ are integral.
Their elementary matrices and inverses therefore lie in $\GL(\OO)$.
The two clearings leave the Schur complement
\begin{equation}\label{eq:eliminationSchur}
 D-ca^{-1}b.
\end{equation}
Every term in it has valuation at least $\alpha$, because
$v(c_i)+v(b_j)-v(a)\ge\alpha$. Cancellation can raise the valuation, but
cannot lower it. Thus subsequent pivot valuations are nondecreasing.
After finitely many steps the result is a diagonal matrix and a zero block.
All transformations were integral and invertible, so they preserve the
$\Delta_k$. The diagonal pivot valuations must therefore be their successive
differences, which are the singular scales by \cref{thm:scales}.
\end{proof}

The algorithm uses $O(mn\min(m,n))$ field operations, apart from pivot
comparisons. This counts algebraic operations; it is not a bit-complexity
bound for arbitrary Hahn-series representations. A
single exact zero test may already be unavailable for a given symbolic
representation. Moreover, the elimination matrices are generally not
unitary: this algorithm returns the singular \emph{scales}, not the singular
values or the best low-rank approximation itself.

\subsection{What must be provided by a symbolic scalar representation}
The algorithm has four logically separate requirements: a representation
closed under the field operations it performs, exact zero testing, exact
leading-exponent access for nonzero outputs, and an effective order on those
exponents. Rational functions over a finite rational or Gaussian-rational
monomial field can meet these requirements. A name for an arbitrary
surreal number does not automatically meet any of them. This is the same
kind of denotation-versus-computation distinction emphasized by the
repository's computer-algebra strand~\cite{RepoMap}.

Cancellation is the principal practical issue. A representation that stores
only the current leading term is insufficient: \eqref{eq:eliminationSchur}
may cancel it. Correct implementations must either compute further exact
coefficients, attach a certified remainder bound, or report that the pivot
scale is unresolved. Reporting a zero pivot merely because all currently
stored coefficients vanish is unsound.

\subsection{Two further contracts: spectral descent and ramification}
\label{spec:sec:contracts}
The support-certified spectral construction of \cref{spec:sec:descent} is
more demanding than the scale algorithm. It needs exact arithmetic and
conjugation, exact zero tests, leading exponents and leading coefficients
for nonzero expressions, ordinary constant Hermitian diagonalization, and a
representation of strongly summable substitutions together with their
support certificates. An exact representation need not provide any of these
effectively merely because it can denote a surreal number.
\Cref{spec:ex:tail} is the concrete obstruction: a procedure that inspects
more and more coefficients never reaches the first nonscalar scale, whereas
the exact trace subtraction \eqref{spec:eq:tracesubtract} reaches it in one
step. A practical implementation should also keep two outputs apart. A
matrix identity such as $U^*AU=D$ is an \emph{algebraic} certificate; a
support certificate is what licenses the Hahn expressions occurring in it.
Finite truncation can test an identity to a declared order, but cannot
certify exact zero, all coefficients, or well-ordering of an arbitrary
support.

The determinantal ramification computation of \cref{spec:sec:ramification}
needs much less: exact valuations of finitely many minors, and a way to
compute the finite subgroup generated by specified classes modulo
$q\Gamma$. In a finite integer lattice \eqref{spec:eq:latticedegree} makes
this a Smith normal form, hence a concrete algorithm.

\begin{warning}[Membership in $q\Gamma$ is extra data]\label{spec:warn:qgamma}
For an abstract ordered group $\Gamma$ presented only through a black-box
comparison of elements, membership of $\Delta_k$ in $q\Gamma$ is additional
group-theoretic data and need not be decidable. Order comparison alone does
not answer it. Consequently
\eqref{spec:eq:criterionmain} --- ``$A^{1/q}$ is defined over $K_\Gamma$ if
and only if every $\Delta_k(A)$ lies in $q\Gamma$'' --- is an exact
mathematical criterion that a given representation may nevertheless be
unable to evaluate, even when it can compute every $\Delta_k$ exactly.
\end{warning}

This is the same denotation-versus-computation distinction as above, one
tier further out, and it is exactly the question that the repository's
computer-algebra report formalizes as a hierarchy of representation
capabilities~\cite{RepoCAS}. That report's own inventory already makes the
matching distinction on the implementation side: a ramified root is
representable in a model whose exponent lattice is $\Q^r$ and is not
representable in one whose exponent lattice is a finite-rank free abelian
group --- which is the
non-divisibility distinction of \cref{spec:sec:ramification} seen from the
other end. \Cref{spec:eq:criterionmain} is therefore a useful concrete
instance of a decision that needs a capability tier the black-box order
interface does not supply. Nothing here claims to solve that report's
representation or computability problems.

\subsection{Why finite truncation cannot certify every exact rank}
Fix any precision threshold $\theta$. Choose $\eta>\theta$. The matrices
\[
 \diag(1,0),\qquad \diag(1,t^\eta)
\]
agree at all exponents at most $\theta$ but have different exact ranks.
Thus finite-precision coefficient access alone cannot decide exact rank
uniformly. An exact representation may still decide it, for example by
computing a nonzero symbolic determinant, but truncation alone cannot.

The right numerical-style output is therefore often a certified rank
profile up to a declared cutoff, with the remaining directions unresolved.
The perturbation theorem in \cref{sec:conditioning} specifies exactly what
precision is sufficient for such a certificate. A finite exact rational
kernel for testing the elimination and Gram identities is included in the
archive; its scope is described in \cref{app:checks}.

\section{Condition numbers and perturbation certificates}\label{sec:conditioning}

\subsection{Distance to singularity is a minimum, not merely an infimum}
This subsection again holds over any real closed $F$ and $K=F(i)$.

\begin{theorem}[Distance to singularity]\label{thm:distance}
For an invertible square $A$,
\begin{equation}\label{eq:dist}
 \norm{A^{-1}}_{\op}=\sigma_n(A)^{-1},\qquad
 \min_{A+E\ \mathrm{singular}}\norm E_{\op}=\sigma_n(A).
\end{equation}
The latter minimum is attained by
$E=-\sigma_nu_nv_n^*$ for any matched least-singular vector pair.
\end{theorem}
\begin{proof}
Invert the SVD to obtain the first formula. The stated perturbation deletes
one diagonal entry of $\Sigma$ and has norm $\sigma_n$. Conversely, if
$A+E$ is singular, choose a unit vector $x$ in its kernel. Then
\[
 \norm E_{\op}\ge\norm{Ex}=\norm{Ax}\ge\sigma_n.
\]
\end{proof}

Define the ordered spectral condition number by
\[
 \kappa(A)=\norm A_{\op}\norm{A^{-1}}_{\op}
          =\frac{\sigma_1}{\sigma_n}\ge1.
\]
It is an element of $F$, not necessarily an ordinary real. In a Hahn
workspace,
\begin{equation}\label{eq:kappaval}
 v(\kappa(A))=\gamma_1-\gamma_n.
\end{equation}
Consequently $\kappa(A)$ is infinite precisely when
$\gamma_1<\gamma_n$. Matrices whose nonzero singular values all have the
same valuation have finite condition number, although its ordinary leading
coefficient can still be large. Overall scale and relative conditioning are
separate: $t^{-\omega}I$ has infinite norm but condition number one.

\begin{theorem}[Inverse and linear-system perturbation]\label{thm:inversepert}
Let $A$ be invertible and put
$q=\norm{A^{-1}}_{\op}\norm E_{\op}$. If $q<1$, then $A+E$ is invertible,
\begin{align}
 \norm{(A+E)^{-1}}_{\op}
 &\le\frac{\norm{A^{-1}}_{\op}}{1-q},\label{eq:invbound}\\
 \norm{(A+E)^{-1}-A^{-1}}_{\op}
 &\le\frac{\norm{A^{-1}}_{\op}^2\norm E_{\op}}{1-q}.
 \label{eq:invdiff}
\end{align}
If $Ax=b\ne0$ and $(A+E)(x+\delta x)=b+f$, then
\begin{equation}\label{eq:forwarderror}
 \frac{\norm{\delta x}}{\norm x}
 \le\frac{\kappa(A)}{1-q}
 \left(\frac{\norm f}{\norm b}
             +\frac{\norm E_{\op}}{\norm A_{\op}}\right).
\end{equation}
\end{theorem}
\begin{proof}
For every vector $x$,
\[
 \norm{(A+E)x}\ge
 (\sigma_n(A)-\norm E_{\op})\norm x.
\]
The coefficient is positive under the hypothesis. The map is injective,
hence invertible in finite dimension, and this inequality gives
\eqref{eq:invbound}. The finite resolvent identity
\[
 (A+E)^{-1}-A^{-1}=-(A+E)^{-1}EA^{-1}
\]
gives \eqref{eq:invdiff}. Also
$\delta x=(A+E)^{-1}(f-Ex)$. Apply \eqref{eq:invbound}, divide by
$\norm x$, and use $\norm b\le\norm A_{\op}\norm x$ to obtain
\eqref{eq:forwarderror}.
\end{proof}

No geometric series was used. This matters because the ordered inequality
$q<1$ is not by itself a Hahn-summability hypothesis.

\subsection{A certified precision theorem for singular scales}
Return to a Hahn workspace. When $v(A)=\eta$, every entry of $A$ has
valuation at least $\eta$; this is the matrix precision convention below.

\begin{theorem}[Preserving visible singular data]\label{thm:precision}
Let $A,B=A+E$ have the same dimensions. If
$v(E)>\gamma_k(A)$ for a nonzero singular value $\sigma_k(A)$, then for
all $1\le j\le k$,
\[
 v(\sigma_j(B))=\gamma_j(A),\qquad
 \lc(\sigma_j(B))=\lc(\sigma_j(A)).
\]
More generally, if $v(E)>\theta$, then
\[
 r_{\le\theta}(A)=r_{\le\theta}(B),\qquad
 r_{<\theta}(A)=r_{<\theta}(B).
\]
The exact ranks need not be equal.
\end{theorem}
\begin{proof}
By \cref{lem:valnorm}, $v(\norm E_{\op})=v(E)$ for $E\ne0$. The
ordered singular-value estimate \eqref{eq:singlip} therefore implies
\[
 v\bigl(\sigma_j(B)-\sigma_j(A)\bigr)\ge v(E).
\]
An error of strictly larger valuation cannot alter a nonzero leading term.
This proves the first assertion. If a singular value has valuation greater
than $\theta$, or is zero, adding an error of valuation greater than
$\theta$ keeps it above that threshold. Singular values at or below the
threshold retain their leading terms. Apply this coordinatewise to the
ordered singular lists for both asserted counts. The case $E=0$ is
immediate.
\end{proof}

Strictness is essential. For
$A=\diag(1,t^\theta)$ with $\theta>0$, a perturbation
$E=\diag(0,-t^\theta)$ has exactly valuation $\theta$ and destroys the
second singular direction. Even under the strict hypothesis, a zero
singular value can become a nonzero value of a still smaller magnitude;
$\diag(1,0)$ and $\diag(1,t^\eta)$ with $\eta>\theta$ show why exact
rank is not certified. The theorem certifies the rank \emph{visible at the
stated scale}.

\section{Nonnormal matrices and ordered pseudospectra}\label{sec:pseudo}

Eigenvalues alone do not control a general matrix. For example,
\[
 N_L=\begin{pmatrix}0&L\\0&0\end{pmatrix}
\]
has only the eigenvalue zero, while $\norm{N_L}_{\op}=|L|$ can be infinite.
The distinction between eigenvalues and singular values is as important in
the surcomplex field as it is over $\C$.

For a square matrix and an ordered tolerance $\varepsilon\in F_{\ge0}$,
define the closed ordered pseudospectrum by
\begin{equation}\label{eq:pseudodef}
 \Lambda_\varepsilon(A)=
 \{z\in K:\sigma_n(zI-A)\le\varepsilon\}.
\end{equation}

\begin{proposition}[Exact perturbation interpretation]\label{prop:pseudo}
A scalar $z$ belongs to $\Lambda_\varepsilon(A)$ if and only if there is
$E$ with $\norm E_{\op}\le\varepsilon$ for which $z$ is an eigenvalue of
$A+E$. If $z\notin\spec A$ and $\varepsilon>0$, this is equivalent to
\[
 \norm{(zI-A)^{-1}}_{\op}\ge\varepsilon^{-1}.
\]
For normal $A$,
\begin{equation}\label{eq:normalpseudo}
 \Lambda_\varepsilon(A)=
   \bigcup_{\lambda\in\spec A}\{z:|z-\lambda|\le\varepsilon\}.
\end{equation}
\end{proposition}
\begin{proof}
Apply the distance-to-singularity argument to $zI-A$. If it is already
singular, use $E=0$. Otherwise \cref{thm:distance} gives an attained
perturbation of norm $\sigma_n(zI-A)$ making it singular; changing its
sign converts it to a perturbation of $A$. The lower bound follows by
applying $zI-A$ to an eigenvector of $A+E$. The inverse characterization is
\eqref{eq:dist}. For normal $A$, unitary diagonalization gives singular
values of $zI-A$ equal to the numbers $|z-\lambda_j|$, proving
\eqref{eq:normalpseudo}.
\end{proof}

\begin{example}[An infinitesimal perturbation with a larger spectral effect]
Let
\[
 N=\begin{pmatrix}0&1\\0&0\end{pmatrix},\qquad
 E=\begin{pmatrix}0&0\\t^2&0\end{pmatrix}.
\]
Then $\norm E_{\op}=t^2$, whereas $N+E$ has eigenvalues $\pm t$.
Thus $t\in\Lambda_{t^2}(N)$ although its distance from the unperturbed
spectrum is $t>t^2$. More generally, for nonzero infinitesimal $z$,
$zI-N$ has $\Delta_1=0$ and $\Delta_2=2v(z)$, so its least singular
value has valuation $2v(z)$. The nonlinear sensitivity is visible directly
in the determinantal scale profile of the resolvent matrix.
\end{example}

Here $\varepsilon$ is an actual ordered magnitude, not just a value-group
element. A valuation threshold gives a coarser notion. In particular,
replacing the ordered disks in \eqref{eq:normalpseudo} by valuation balls
would change the assertion. No fine-continuous contour is needed in the
pseudospectral argument.

\section{Spectral subspaces and infinitesimal gaps}\label{sec:subspaces}

Individual eigenvectors are not canonical, especially at multiplicities.
Orthogonal spectral projectors are the better perturbation objects. The
following finite-dimensional Frobenius version of a Davis--Kahan bound is
proved directly over $F$; the classical context is
\cite{DK,YWS}. No inverse trigonometric function is needed to measure the
subspace displacement.

\begin{theorem}[A separated-subspace bound]\label{thm:subspace}
Let $H=H^*$, $B=H+E=B^*$, and let $P$ be an orthogonal projector onto an
$H$-invariant subspace. Let $Q$ project onto a $B$-invariant subspace of
rank $k$. Suppose that every eigenvalue $\lambda$ of
$H|_{\im(I-P)}$ and every eigenvalue $\mu$ of $B|_{\im Q}$ satisfy
\[
 |\lambda-\mu|\ge\delta>0.
\]
Then
\begin{equation}\label{eq:subspacebound}
 \norm{(I-P)Q}_{\Fr}\le\frac{\norm E_{\Fr}}{\delta}.
\end{equation}
If $\rank P=\rank Q$, also
\begin{equation}\label{eq:projbound}
 \norm{P-Q}_{\Fr}\le\frac{\sqrt2\norm E_{\Fr}}{\delta}.
\end{equation}
\end{theorem}
\begin{proof}
Choose orthonormal eigenbasis matrices $Y$ for $\im(I-P)$ and $X$ for
$\im Q$, so $HY=Y\Lambda$ and $BX=XM$, with real diagonal
$\Lambda,M$. Put $C=Y^*X$. Multiplying $HX+EX=XM$ on the left by $Y^*$
gives the Sylvester equation
\[
 \Lambda C-CM=-Y^*EX.
\]
Entrywise,
$(\lambda_a-\mu_b)c_{ab}=-(Y^*EX)_{ab}$, hence
\[
 \norm C_{\Fr}^2\le\delta^{-2}\norm{Y^*EX}_{\Fr}^2
                         \le\delta^{-2}\norm E_{\Fr}^2.
\]
The last inequality follows because $X$ and $Y$ have orthonormal columns;
extend them to unitary matrices and observe that taking a submatrix cannot
increase the sum of squared moduli. Since $(I-P)Q=YCX^*$,
$\norm{(I-P)Q}_{\Fr}=\norm C_{\Fr}$. This proves
\eqref{eq:subspacebound}.

For equal ranks,
\[
 \norm{P-Q}_{\Fr}^2
 =\tr P+\tr Q-2\tr(PQ)
 =2\norm{(I-P)Q}_{\Fr}^2,
\]
which gives \eqref{eq:projbound}.
\end{proof}

\begin{corollary}[Using an unperturbed cluster gap]\label{cor:clustergap}
Select a nonempty proper set of eigenvalue indices of $H$, separated from
all complementary indices by at least $g>0$. Let $P$ project onto this
cluster, and let $Q$ select the same indices in the decreasing eigenvalue
list of $H+E$. If $e=\norm E_{\op}<g/2$, then
\[
 \norm{P-Q}_{\Fr}\le\frac{\sqrt2\norm E_{\Fr}}{g-e}.
\]
\end{corollary}
\begin{proof}
The eigenvalue estimate in \cref{thm:weyl} moves each selected eigenvalue
by at most $e$. Its distance from an unselected eigenvalue of $H$ is at
least $g-e$. No eigenvalue can cross the external gap under $e<g/2$.
The chosen subspaces have the same dimension, so apply
\cref{thm:subspace} with $\delta=g-e$.
\end{proof}

The Frobenius separation condition is all that was used in the proof.
We are not asserting every operator-norm variant of the classical
Davis--Kahan theorem under an arbitrary separation hypothesis. The
Frobenius estimate is sufficient for an important surreal consequence:
if $v(E)>v(g)$ in a Hahn workspace, then the projector displacement is
infinitesimal. Indeed, $\norm E_{\Fr}/g$ then has positive valuation, so
$e<g/2$ and the displayed bound has positive valuation.

\begin{example}[Infinitesimal error, ordinary subspace rotation]\label{ex:rotation}
Consider
\[
 H=\begin{pmatrix}0&0\\0&t^2\end{pmatrix},\qquad
 B=\frac{t^2}{2}\begin{pmatrix}1&-1\\-1&1\end{pmatrix}.
\]
They have the same eigenvalues $0,t^2$. Their zero-eigenspace projectors are
\[
 P=\begin{pmatrix}1&0\\0&0\end{pmatrix},\qquad
 Q=\frac12\begin{pmatrix}1&1\\1&1\end{pmatrix}.
\]
Exact multiplication gives
\[
 \norm{B-H}_{\op}=\frac{t^2}{\sqrt2},\qquad
 \norm{P-Q}_{\Fr}=1,\qquad \norm{P-Q}_{\op}=\frac1{\sqrt2}.
\]
Thus the matrix error is infinitesimal but the subspace change is not.
Its size is comparable to the gap $t^2$, so the gap-relative smallness
hypothesis fails. Discarding the common infinitesimal factor before asking
about eigenvectors would conceal the entire phenomenon.
\end{example}

\subsection{Singular subspaces and a warning about Gram squaring}
Right singular subspaces are eigenspaces of $A^*A$. Under $A\mapsto A+E$,
\[
 (A+E)^*(A+E)-A^*A=A^*E+E^*A+E^*E,
\]
so its operator norm is at most
$2\norm A_{\op}\norm E_{\op}+\norm E_{\op}^2$.
One can therefore apply \cref{thm:subspace} to squared-singular-value gaps.
This is valid but can be quantitatively wasteful: it changes gaps
$\sigma_i-\sigma_j$ into $\sigma_i^2-\sigma_j^2$. A direct singular-subspace
or Hermitian-dilation argument can yield better bounds, but no sharp Wedin
theorem is claimed here. The distinction is particularly relevant when
both singular values are infinitesimal.

\section{Degenerate clusters and Schur-complement corrections}\label{sec:schur}

A cluster of repeated or nearly repeated eigenvalues cannot generally be
understood by retaining only the perturbation inside that cluster. Coupling
to its complement can create a smaller, second-order splitting. The right
starting point is an exact finite identity.

\begin{proposition}[Exact effective matrix]\label{prop:effective}
Suppose a Hermitian matrix is written in an orthogonal decomposition as
\[
 H=\begin{pmatrix}
       \lambda_0 I+E_{11}&E_{12}\\
       E_{21}&D+E_{22}
    \end{pmatrix},\qquad E_{21}=E_{12}^*.
\]
For $z\in K$ such that $D+E_{22}-zI$ is invertible,
\begin{align}
 \det(H-zI)&=\det(D+E_{22}-zI)\det S(z),\label{eq:Schurdet}\\
 S(z)&=(\lambda_0-z)I+E_{11}
        -E_{12}(D+E_{22}-zI)^{-1}E_{21}.\label{eq:effective}
\end{align}
If the ordered distance from $z$ to every eigenvalue of $D+E_{22}$ is at
least $\delta>0$, then
\begin{equation}\label{eq:effectivebound}
 \norm{E_{12}(D+E_{22}-zI)^{-1}E_{21}}_{\op}
 \le\frac{\norm{E_{12}}_{\op}^2}{\delta}.
\end{equation}
\end{proposition}
\begin{proof}
Multiply $H-zI$ on the left by the determinant-one block matrix
\[
 \begin{pmatrix}I&-E_{12}(D+E_{22}-zI)^{-1}\\0&I\end{pmatrix}.
\]
The resulting matrix is block lower triangular with diagonal blocks $S(z)$
and $D+E_{22}-zI$, proving \eqref{eq:Schurdet}. Diagonalization of the
Hermitian complement shows that its shifted inverse has norm at most
$\delta^{-1}$. Apply the product norm bounds and
$\norm{E_{21}}_{\op}=\norm{E_{12}}_{\op}$.
\end{proof}

For real $z$ outside the complementary spectrum, $S(z)$ is Hermitian. It is
an eigenvalue-dependent effective matrix, not in general one fixed matrix
whose eigenvalues already solve the original problem. Freezing $z$ or
dropping the inverse correction requires a separate error estimate.
Formula \eqref{eq:effectivebound} supplies one such estimate: the correction
has the scale of the square of the coupling divided by the complementary
gap. If that gap is itself infinitesimal, a seemingly second-order term can
be as important as a first-order term.

\begin{example}[A spectral shift invisible in the small diagonal block]\label{ex:secondorder}
Let
\[
 H_t=\begin{pmatrix}0&t\\t&1\end{pmatrix}.
\]
Its eigenvalues and low-eigenvalue projector are exactly
\begin{gather}
 \lambda_\pm=\frac{1\pm\sqrt{1+4t^2}}2,\label{eq:secondroots}\\
 P_- =\frac{\lambda_+ I-H_t}{\lambda_+-\lambda_-}.\label{eq:secondprojector}
\end{gather}
The small eigenvalue is
\begin{equation}\label{eq:secondexpansion}
 \lambda_-=-t^2+t^4-2t^6+5t^8+O_v(t^{10}).
\end{equation}
Here $O_v(t^{10})$ denotes a Hahn remainder of valuation at least $10$;
$t$ is not a real variable tending to zero. The expansion follows either
from the binomial series justified in \cref{sec:functions}, or by substituting
the displayed finite polynomial into $\lambda^2-\lambda-t^2$ and checking
the remainder order, then selecting the root with positive valuation.

The direct perturbation inside the first coordinate is zero. Nevertheless,
\eqref{eq:effective} gives the scalar equation
$-z-t^2/(1-z)=0$, whose leading solution is $z=-t^2$.
An eigenvector may be taken as $(1,\lambda_-/t)^{\mathsf T}$, displaying
an order-$t$ mixing with a spectral shift of order $t^2$.
\end{example}

The companion matrix
\[
 \widetilde H_t=\begin{pmatrix}t^2&t\\t&1\end{pmatrix}
 =\begin{pmatrix}t\\1\end{pmatrix}\begin{pmatrix}t&1\end{pmatrix}
\]
is positive semidefinite of exact rank one. In its effective scalar at
$z=0$, the direct $t^2$ term and the coupling correction $t^2$ cancel
exactly. The difference between a small eigenvalue and a zero eigenvalue
can therefore depend on cancellation between terms of the same scale.
Neither leading entry valuations nor a generic ``second-order is negligible''
rule resolves it.

\section{Finite functional calculus versus infinite Hahn series}\label{sec:functions}

\subsection{Finite spectral calculus needs no convergence}
Let $H$ be Hermitian with distinct eigenvalues in a finite set $S$. For
any prescribed scalars $f(\lambda)\in K$, define
\begin{equation}\label{eq:finitecalc}
 f(H)=\sum_{\lambda\in S}f(\lambda)P_\lambda.
\end{equation}
Lagrange interpolation gives a polynomial $p\in K[X]$ of degree less than
$|S|$ with $p(\lambda)=f(\lambda)$; thus $f(H)=p(H)$. The construction
preserves finite sums and products and depends only on the values of $f$
at the spectrum. For real-valued spectral data it gives a Hermitian matrix;
for nonnegative data it gives a positive semidefinite matrix. It also yields
\[
 \norm{f(H)}_{\op}=\max_{\lambda\in S}|f(\lambda)|.
\]
For $z\notin S$, the resolvent is
\[
 (zI-H)^{-1}=\sum_{\lambda\in S}\frac{P_\lambda}{z-\lambda}.
\]
All of this is finite algebra. It does not assert that an arbitrary scalar
function already has a chosen surreal extension. In particular, one cannot
use \eqref{eq:finitecalc} to define unspecified values such as a global
surreal sine at an infinite argument by circular reasoning. One must first
supply the values to be interpolated.

\subsection{A genuine matrix Hahn-summability theorem}
A family of Hahn series is \emph{strongly summable} if the union of its
supports is well ordered and every exponent occurs in only finitely many
members (the notion of \cref{spec:def:summable}). Its sum is then defined
coefficientwise by finite sums. For matrices this means entrywise strong
summability, equivalently the same condition on the finitely many entry
families together.

We use the standard support lemma underlying Hahn-field arithmetic: if
$S\subseteq\Gamma_{>0}$ is well ordered, its additive monoid
\[
 S^*=\{0\}\cup\{s_1+\cdots+s_k:k\ge1,\ s_j\in S\}
\]
is well ordered, and every exponent has only finitely many representations
as an ordered finite sum of members of $S$. This is the positive-support
form of the Neumann lemma; see the support development in Gonshor,
Chapter~5A~\cite{Gonshor}, and \cref{spec:lem:neumann} for a proof from the
ordered-word lemma. The lemma does not assert that $k\rho$ is cofinal in
$\Gamma$ for a positive $\rho$.

\begin{theorem}[Matrix inverse and square-root series]\label{thm:series}
If $E\in K_\Gamma^{n\times n}$ has $v(E)>0$, then
\begin{equation}\label{eq:Neumann}
 \sum_{k\ge0}E^k
\end{equation}
is strongly Hahn-summable and equals $(I-E)^{-1}$. If in addition $E=E^*$,
then
\begin{equation}\label{eq:binomial}
 \sum_{k\ge0}\binom{1/2}{k}E^k
\end{equation}
is strongly Hahn-summable and equals the positive square root of $I+E$.
\end{theorem}
\begin{proof}
The finite union of the entry supports of $E$ is a well-ordered subset
$S\subseteq\Gamma_{>0}$. An entry of $E^k$ is a finite sum of products
of $k$ entries, so its support is contained in sums of $k$ elements of $S$.
The support lemma gives a well-ordered union and finitely many contributing
lengths and tuples at every exponent. For a fixed length there are only
finitely many matrix-index paths. This proves strong summability of both
series, and justifies the coefficientwise rearrangements below.

The finite formal identity $(1-X)\sum_{k\ge0}X^k=1$ is valid
coefficientwise after substituting $E$, giving a two-sided inverse. Likewise
the formal binomial identity makes the square of \eqref{eq:binomial} equal
to $I+E$. Its sum $R$ is Hermitian, because every power of $E$ is, and
$v(R-I)>0$. The norm of $R-I$ is therefore infinitesimal by
\cref{lem:valnorm}; all eigenvalues of $R$ are positive by
\cref{thm:weyl}. Uniqueness of the positive square root from
\cref{prop:positive} identifies $R$ with $(I+E)^{1/2}$.
\end{proof}

\subsection{Two failures of a familiar shortcut}
The ordered condition $\norm E_{\op}<1$ does not imply strong
Hahn-summability of \eqref{eq:Neumann}. For $E=\tfrac12 I$, every nonzero
entry of every power has exponent zero, so infinitely many members meet
the same exponent. The Hahn summability criterion fails, although
$(I-E)^{-1}=2I$ exists algebraically. Ordinary analytic summation of real
coefficients is a different operation.

Even when the Hahn series exists, it need not be a topological limit of its
partial sums in the full surreal field. With $E=tI$ and $t=\omega^{-1}$,
\[
 (I-tI)^{-1}-\sum_{k=0}^{N}t^kI=\frac{t^{N+1}}{1-t}I.
\]
For every ordinary $N$, this error is larger in modulus than $t^\omega$.
Thus the partial sums fail the fine-surreal convergence test at the single
tolerance $t^\omega$.

There is no contradiction with internal valuation convergence in
$K_\Q=\C((t^\Q))$: the natural-number exponents are cofinal in $\Q$, so
these same partial sums do converge in that field's own valuation topology.
In the larger workspace with value group $\Q+\Q\omega$, they do not,
because the threshold $\omega$ is now present. Algebraic identities,
strong summability, and workspace-relative topology are three different
claims; only the first two were needed for \cref{thm:series}.

\section{Two applications: regularization and finite tensors}\label{sec:applications}

\subsection{Regularization as an explicit scale filter}
For $\alpha>0$, consider the regularized least-squares objective
\[
 J_\alpha(x)=\norm{Ax-b}^2+\alpha^2\norm x^2.
\]
The matrix $A^*A+\alpha^2I$ is positive definite, regardless of the rank of
$A$. Define
\begin{equation}\label{eq:tikhonov}
 x_\alpha=(A^*A+\alpha^2I)^{-1}A^*b.
\end{equation}
Expanding and using its normal equation yields
\[
 J_\alpha(x)=J_\alpha(x_\alpha)
 +(x-x_\alpha)^*(A^*A+\alpha^2I)(x-x_\alpha).
\]
Hence \eqref{eq:tikhonov} is the unique minimizer over the whole field, with
no compactness or coercivity argument required.

In singular coordinates,
\begin{equation}\label{eq:filters}
 x_\alpha=\sum_{j=1}^r
    \frac{\sigma_j}{\sigma_j^2+\alpha^2}(u_j^*b)v_j,
 \qquad
 Ax_\alpha=\sum_{j=1}^r
    \frac{\sigma_j^2}{\sigma_j^2+\alpha^2}(u_j^*b)u_j.
\end{equation}
The fitted-output factor has a particularly clear reduction. Put
$\beta=v(\alpha)$ and $\gamma=v(\sigma_j)$, with positive leading
coefficients $b_0=\lc(\alpha)$ and $a_0=\lc(\sigma_j)$. Positivity of
the denominator gives
\begin{equation}\label{eq:regscale}
 \st\left(\frac{\sigma_j^2}{\sigma_j^2+\alpha^2}\right)
 =\begin{cases}
  1,&\gamma<\beta,\\
  \displaystyle\frac{a_0^2}{a_0^2+b_0^2},&\gamma=\beta,\\
  0,&\gamma>\beta.
 \end{cases}
\end{equation}
The quotient always lies in $[0,1]$, so its standard part exists. This is a
scale filter: larger singular directions pass, smaller ones disappear in
the ordinary shadow, and the boundary scale retains a nontrivial leading
amplitude. At that boundary, a valuation-only description is insufficient.

\begin{proposition}[Ordinary data and infinitesimal regularization]
If $A$ and $b$ have ordinary complex entries and $\alpha$ is a positive
infinitesimal in a Hahn workspace, then $x_\alpha$ is finite and
\[
 \st(x_\alpha)=A^\dagger b,
\]
where the right side is the ordinary complex pseudoinverse solution.
\end{proposition}
\begin{proof}
Choose an ordinary complex SVD. Every nonzero singular value is a positive
ordinary real number. In \eqref{eq:filters}, the coefficient
$\sigma_j/(\sigma_j^2+\alpha^2)$ is finite and has standard part
$1/\sigma_j$. There are finitely many terms, so reduction gives the
pseudoinverse formula.
\end{proof}

The ordinary-data hypothesis matters. For $A=(t)$, $b=1$, and
$\alpha=t^2$, the regularized solution is
$t^{-1}/(1+t^2)$, which is infinite and has no standard part. Nor does a
universal choice of one infinitesimal $\alpha$ resolve every surreal scale:
\cref{ex:three} already exhibits a direction smaller than all finite powers
of a chosen basic infinitesimal.

\subsection{Schmidt decomposition and infinitesimal components}
Give the finite tensor space $K^m\otimes_K K^n$ the inner product for which
the coordinate tensor basis is orthonormal. A vector
\[
 \psi=\sum_{a,b}A_{ab}e_a\otimes f_b
\]
is identified with its coefficient matrix $A$. The SVD immediately gives
\begin{equation}\label{eq:Schmidt}
 \psi=\sum_{j=1}^r\sigma_j u_j\otimes\bar v_j,
 \qquad \norm\psi^2=\sum_{j=1}^r\sigma_j^2.
\end{equation}
Conjugation on the second vector is required by the convention
$A=U\Sigma V^*$. The least number of simple tensors in a representation
of $\psi$ is its matrix rank: a sum of $k$ simple tensors gives a matrix
of rank at most $k$, and \eqref{eq:Schmidt} attains $k=r$.

For a unit vector, the eigenvalues of the positive trace-one matrix $AA^*$
are the $\sigma_j^2$ and zeros. Thus the finite algebraic parts of Schmidt
rank and reduced-matrix spectral calculations make sense over $K$.
For example,
\[
 \psi_t=\frac{e_0\otimes f_0+t\,e_1\otimes f_1}{\sqrt{1+t^2}}
\]
has exact tensor rank two, whereas its standard part is the rank-one tensor
$e_0\otimes f_0$. Its second squared Schmidt coefficient is the positive
infinitesimal $t^2/(1+t^2)$, not zero. This illustrates precisely what is
lost on taking an ordinary shadow.

This is a finite algebraic application, not a proposed physical
interpretation of surreal probabilities. No infinite tensor product,
measure-valued probability theory, entropy extension, or quantum dynamics
is supplied by these calculations.

\section{What the chapter establishes, and what remains separate}\label{sec:scope}

The finite-dimensional order theory is robust: orthogonality, spectral
and singular-value decompositions, least squares, low-rank optimization,
and the perturbation bounds proved above require real closedness, not
Archimedean completeness. Their proofs exhibit the relevant witnesses, so
questions about class-sized compactness never arise.

The scale theory contains additional information not visible in an ordinary
complex matrix. Nonzero singular values have a finite ordered list of
valuations. Determinantal minima recover that list, positive Gram
coefficients recover its leading amplitudes in blocks, and the input
filtration gives an intrinsic meaning to rank at a chosen resolution.
The same data govern conditioning and perturbation certificates. The
minimum-valuation elimination gives a finite algebraic route to those
invariants when the scalar representation supplies the required exact
operations.

Divisibility of the value group turns out to be separable from all of this.
The decompositions survive its removal, with every exponent confined to the
subgroup generated by the input exponents and with a finite splitting tree
of at most $n-1$ nontrivial nodes, a bound independent of the order type of
the supports and of the rank or cardinality of $\Gamma$. What does not
survive its removal is the free availability of positive roots, and the
determinantal invariant measures the deficit exactly: the field of
definition of $A^{1/q}$, its degree, the criterion for no extension at all,
the matching criterion for a Gram factorization, and a single scalar
generator of the whole extension. That the same $\Delta_k$ answers both the
scale question and the arithmetic one is the point at which the two halves
of this article meet.

Several nearby subjects remain distinct. There is no theorem here about
infinite-dimensional Hilbert spaces over $\SC$, bounded-operator spectra,
compact or trace-class operators, or spectral measures. A support-controlled
Hahn expansion of an eigenbasis \emph{is} now available for a single matrix,
over an arbitrary set-sized value group
(\cref{spec:thm:hermitian,spec:cor:normal,spec:thm:svd}). There is still no
global analytic choice of eigenvectors through collisions or across multiple
parameters, and \cref{spec:warn:crossing} shows that this is not a gap in
the argument but a genuine failure in the multivariable formal
power-series category: the crossing matrix \eqref{spec:eq:crossing} has no
eigenvalue in $\C[[x,y]]$ at all. The descent theorems resolve the totally
ordered scale problem and not that one. Nor do the exact decompositions
automatically provide effective algorithms for arbitrary surreal inputs;
\cref{spec:sec:contracts} lists what a representation must supply, and
\cref{spec:warn:qgamma} records that even the ramification criterion can
outrun a black-box order interface. Finally, the strong matrix series of
\cref{thm:series} is not a claim of fine-topological convergence, and
neither is the splitting tree: \cref{spec:ex:nonconvergence} exhibits a
value group in which a valid Hahn identity has non-convergent partial sums,
so a remainder bound is never used here as a convergence proof.

All field and extension statements here are set-sized. No proper-class Hahn
sum, no proper-class vector space, and no class-sized workspace occurs
anywhere; \cref{spec:prop:exactlocalization} is what keeps the surcomplex
statements inside a set.

Three further questions are concrete enough to name. One is to replace the
finite tree of set-valued supports by finite effective support grammars for
useful transseries subclasses. Another is to characterize smaller non-Hahn
subfields closed under the separated-block construction, so that the full
Hahn field can be replaced by a more economical domain. A third is to ask
which other positive spectral functionals, beyond rational powers and their
strictly positive sums, furnish similarly exact primitive generators. These
are questions suggested by the results, not advertised as previously
published open problems.

A formalization would have a natural order, which is worth recording
because it is \emph{not} a report of code that has been compiled: positive
sum-of-squares valuation (\cref{spec:lem:normroot}) and integral unitary
invariance (\cref{lem:unitaryintegral}) first; then finite-index Hahn coset
decompositions, the monomial character action, and the positive-coset
primitive-element argument (\cref{spec:lem:finiteindex,spec:lem:positivetrace}),
all of which are set-sized and independent of the full surreal construction;
then an interface for finite-variable formal series evaluated at Hahn
infinitesimals, with the universal separated-block projectors and their
polar intertwiner, where the ordinary complex contour justification in
\cref{spec:thm:split} could be replaced by formal resolvent identities or a
formal Hensel block-lifting theorem; then dimension induction for the
splitting tree; and finally the normal-form embeddings. Nothing in this
article has been proof-assistant verified.

For the repository, this report sits at
\texttt{docs/surcomplex/spectral-theory/}. It follows the workspace
conventions of the foundations and polynomial-algebra reports and can be
read independently of the analytic coefficient-ring machinery. The most
useful cross-reference from polynomial algebra is the bridge from its Newton
profiles to \cref{thm:gram,thm:residual}; the most useful cross-references
to computer algebra are the exact-operation contract in \cref{sec:algorithm}
and its two extensions in \cref{spec:sec:contracts}, where
\cref{spec:warn:qgamma} names a decision that needs a capability tier the
black-box order interface does not provide~\cite{RepoCAS}.

\appendix
\section{A proof and hypothesis ledger}\label{app:ledger}

\begin{longtable}{@{}P{.30\textwidth}P{.28\textwidth}P{.34\textwidth}@{}}
\toprule
Result & Scalar hypotheses & Essential finite mechanism\\
\midrule\endfirsthead
\toprule Result & Scalar hypotheses & Essential finite mechanism\\\midrule\endhead
Spectral theorem and SVD & Real closed $F$, $K=F(i)$ & Algebraic roots,
positivity, finite orthogonal complements\\
Variational and low-rank extrema & Same & Eigenbasis witnesses and dimension
intersection, not compactness\\
Least squares and pseudoinverse & Same & Orthogonal coordinates and finite
sums of nonnegative squares\\
Singular scales from minors & Divisible Hahn workspace & Integral unitary
factors and Cauchy--Binet\\
Leading singular amplitudes & Same, real/complex residue field & Positive
sum of squared minor coefficients\\
Rank filtration & Same, chosen integral lattice & Lattice preservation by
unitaries and coordinate reduction\\
Scale elimination & Exact field and valuation operations & Minimum pivot,
integral operations, Schur complement\\
Inverse perturbation & Real closed $F$; $q<1$ & A positive lower singular
bound and a finite inverse identity\\
Visible-rank stability & Hahn workspace; strict precision threshold & Ordered
singular-value Lipschitz bound plus leading-term preservation\\
Spectral-subspace bound & Real closed $F$; separated selected/complementary
spectra & Entrywise Sylvester equation and Frobenius norm\\
Infinite matrix series & Hahn workspace; $v(E)>0$ & Positive support lemma,
not cofinality of integer multiples\\
\addlinespace
Positive norm square without real closure & Real coefficients, leading
order, positive-support binomial evaluation & Square roots of \emph{all}
positive field elements are not assumed\\
Separated-block lift & Distinct ordinary residue eigenvalues,
finite-variable formal identities, positive Hahn support & Not sequential
valuation convergence, not a surreal contour integral\\
Hermitian spectral descent & Exact trace subtraction, ordinary constant
spectral theorem, block lift, induction on dimension & Divisibility of
$\Gamma$, real closedness of $F_\Gamma$, finite rank\\
Normal and simultaneous descent & Hermitian descent, commutant argument,
finite-dimensional refinement & A pre-existing algebraic closure of
$K_\Gamma$, or an infinite product of unitaries\\
SVD without divisibility & Hermitian Gram diagonalization, norm-square
roots, finite Gram--Schmidt & Square roots of arbitrary positive eigenvalues
of arbitrary positive matrices\\
Exact root field & Scalar binomial roots, spectral descent, finite-index
coset basis & A generic-eigenvalue assumption, nonzero determinant, or
simple spectrum\\
Primitive trace & Roots of unity, coset character action, \emph{strictly
positive} coefficients & Generic weights, independent eigenvalue classes, or
an arbitrary square-root branch\\
Surcomplex interpretation & Set-supported normal forms, set-sized input
subgroup & Proper-class sums or vector spaces\\
\bottomrule
\end{longtable}

The order-theoretic arguments also apply directly to real symmetric
matrices over $F$. They do not apply to an arbitrary algebraically closed
field equipped with an arbitrary involution: positivity of $x^*x$ is the
indispensable additional structure. Conversely, the valuation-ring
elimination argument itself requires no conjugation or positivity; those
structures are what identify its invariant factors with singular values.
This separation helps avoid attributing a positivity theorem to general
valued-field linear algebra.

\section{Reproducible finite checks}\label{app:checks}

The archive contains \texttt{code/verify\_examples.py}, an exact SymPy
verification program, together with its recorded JSON and text reports.
The valuation and elimination kernel uses rational functions over $\Q(i)$
in two formal positive monomials
$t,u$, with the lexicographic valuation convention
\[
 v(t)=(0,1),\qquad v(u)=(1,0).
\]
These correspond to the exponents $1$ and $\omega$ in the subgroup
$\Z+\Z\omega$. A rational function is first canceled exactly; the
valuation is the least lexicographic monomial exponent of its numerator
minus that of its denominator. This finite rational-function field embeds
in a divisible Hahn workspace, but the code is not a general Hahn-series
engine. Separate finite identity checks also use exact square-root
expressions for the two-by-two spectral examples.

The checks compare minimum-valuation elimination with exhaustive minors on
small exact matrices, including matrices obtained by integral invertible
row and column operations from prescribed diagonal scales. They also check
the positive Gram coefficient identity, the residual polynomial examples,
Penrose identities for a rank-deficient matrix, the exact projector and
rotation examples, the second-order spectral expansion, the nilpotent
pseudospectral example, and the regularization and tensor identities.
All symbolic assertions must pass for the program to write a successful
report; no floating-point sampling is used.

The bundled run passed 176 exact assertions across 18 matrix-profile cases
(Python 3.13.5, SymPy 1.14.0). The recorded run is a reproducibility check of these finite identities and
of the implementation on its stated test cases. It is not verification of
the general mathematical proofs, of arbitrary ordered groups or Hahn
supports, or of proper-class foundations. In particular, a finite test
cannot prove the support lemma or certify a general surreal equality
algorithm. The exact assertion count and software versions are recorded in
\texttt{data/verification\_report.json} rather than being used as a proxy
for mathematical correctness.

\subsection*{The second check suite}
The material of
\cref{spec:sec:descent,spec:sec:ramification,spec:sec:trace} arrived with
its own exact SymPy program, also with no floating-point sampling. It
exercises the resolvent-projector construction of
\eqref{spec:eq:projresolvent} and its polar intertwiner
\eqref{spec:eq:polarintertwiner} through formal degree six for a
three-dimensional Hermitian example with a repeated leading eigenvalue; it
tests the displayed scalar and two-by-two identities; it checks minor-valued
profiles under rational unitary changes of basis, including rank-deficient
profiles; and it tests finite quotient computations, subgroup enumeration
against a maximal-minor lattice index, and character and stabilizer counts
in finite integer lattices, the last of these being the computational face
of \eqref{spec:eq:latticedegree} and \eqref{spec:eq:orbitpoly}.

The delivered run passed \textbf{807 exact scalar assertions} --- 378
projector and polar-series coefficients, 144 finite ramification-group
checks, 105 rational-matrix checks, 96 character checks, 55 block and
low-degree terms, 29 displayed identities --- across 48 finite quotient
cases and six rational positive semidefinite profiles, under Python 3.13.5
and SymPy 1.14.0. It was rerun for this collection with an explicit output
directory and reproduced \textsc{pass} with exactly 807 assertions and the
identical category split under Python 3.14.4 and SymPy 1.14.0.

The same caveats apply, and one more. The counts include individual zero
coefficients, so they are not counts of independent theorems. The checks are
finite identities and examples only: there is no exhaustive computation over
arbitrary Hahn supports or arbitrary ordered groups, and the general support
and field-theoretic results rest on the proofs in this article and not on
the scripts. Both programs default their output directory to the shipped
\texttt{data/}, so both must be run on a copy or with an explicit output
path; the repository \textsc{readme} records the safe invocations.

\section{Notation at a glance}
\begin{longtable}{@{}P{.27\textwidth}P{.67\textwidth}@{}}
\toprule Symbol & Meaning\\\midrule\endfirsthead
\toprule Symbol & Meaning\\\midrule\endhead
$F$, $K=F(i)$ & A real closed ordered field and its algebraic complexification\\
$F_\Gamma$, $K_\Gamma$ & Real and complex Hahn workspaces, $t^\gamma=\omega^{-\gamma}$\\
$A^*$ & Conjugate transpose\\
$\sigma_1\ge\cdots\ge\sigma_r>0$ & Nonzero singular values, in decreasing ordered magnitude\\
$\norm A_{\op}$, $\norm A_{\Fr}$ & Operator and Frobenius norms, valued in $F_{\ge0}$\\
$v(a)$, $\lc(a)$ & Least Hahn exponent and coefficient at that exponent\\
$v(A)$ & Minimum valuation of the entries; $v(0)=\infty$\\
$\gamma_j$ & $v(\sigma_j)$, in nondecreasing order\\
$\Delta_k$ & Minimum valuation of a $k$-minor, with $\Delta_0=0$\\
$\OO$, $\mm$, $\st$ & Integral valuation ring, positive-valuation ideal, residue map\\
$c_k$ & Coefficient of $s^k$ in $\det(I+sA^*A)$\\
$C_k$ & Positive leading coefficient of $c_k$, after monomial normalization\\
$r_{<\theta}$, $r_{\le\theta}$ & Numbers of nonzero singular scales below or at a cutoff\\
$A^\dagger$ & Moore--Penrose inverse\\
$\kappa(A)$ & Ordered spectral condition number of an invertible matrix\\
$O_v(t^\eta)$ & A formal remainder with valuation at least $\eta$, not a real limit\\
$\Gamma_\Q$ & The divisible hull $\Gamma\otimes_\Z\Q$, with its unique compatible order\\
$\langle S\rangle_\N$ & The additive monoid generated by a well-ordered $S\subseteq\Gamma_{>0}$\\
$H$ & The exact input support subgroup of a matrix, not assumed divisible\\
$\Gamma_{A,q}$ & The root group $\Gamma+\sum_k\Z\Delta_k(A)/q$\\
$R_q(A)$ & The subgroup of $\Gamma/q\Gamma$ generated by the classes $\Delta_k(A)+q\Gamma$\\
$\mu_q$ & The group of $q$th roots of unity in $\C^\times$\\
$K(B)$ & The field obtained from $K$ by adjoining every entry of a finite matrix $B$\\
$O_{\mathrm{formal}}(E^k)$ & A support-certified sum of monomials of formal degree at least $k$\\
\bottomrule
\end{longtable}

\begin{thebibliography}{99}
\raggedright
\addcontentsline{toc}{section}{References}

\bibitem{RepoMap}
Vladimir Reshetnikov, \emph{The surreal and surcomplex reports}, repository
catalogue, \texttt{docs/README.md}, revision
\texttt{\repocommit}, inspected September 21, 2026.
\href{https://github.com/VladimirReshetnikov/Surreal/blob/e260237db9b71da8b74a0c13c8e6355119091100/docs/README.md}{Pinned catalogue}.
AI-assisted repository drafts; used as documentation-coverage evidence.

\bibitem{RepoPoly}
\emph{Surcomplex Polynomial Algebra: Root Geometry, Newton Profiles,
Multiscale Stability, and Residue Duality}, repository report and detailed
inventory, \texttt{docs/surcomplex/polynomial-algebra/}, same revision as
\cite{RepoMap}. The inventory and the opening source sections were inspected.
\href{https://github.com/VladimirReshetnikov/Surreal/blob/e260237db9b71da8b74a0c13c8e6355119091100/docs/surcomplex/polynomial-algebra/README.md}{Pinned inventory}.
Used to identify adjacent coverage, not as a substitute for the proofs here.

\bibitem{RepoFound}
\emph{Foundations for surreal and surcomplex mathematics}, repository
inventory, \texttt{docs/foundations-and-computation/foundations/README.md},
same revision as \cite{RepoMap}.
\href{https://github.com/VladimirReshetnikov/Surreal/blob/e260237db9b71da8b74a0c13c8e6355119091100/docs/foundations-and-computation/foundations/README.md}{Pinned foundations inventory}.

\bibitem{RepoCAS}
\emph{Computer algebra}, repository report,
\texttt{docs/foundations-and-computation/computer-algebra/}, same revision
as \cite{RepoMap}. Its hierarchy of representation capabilities --- exact
denotation, coefficient access, decidable equality, approximation --- is the
setting for the exact-operation contracts of
\cref{sec:algorithm,spec:sec:contracts}. Consulted through the repository's
documentation map, not through a full audit of its implementations.

\bibitem{Gonshor}
Harry Gonshor, \emph{An Introduction to the Theory of Surreal Numbers},
London Mathematical Society Lecture Note Series 110, Cambridge University
Press, 1986. In particular Chapter~5, on the semigroup support lemma,
normal forms, and real closure.
\href{https://www.cambridge.org/core/books/an-introduction-to-the-theory-of-surreal-numbers/312AE504A3E88E804054BFB390446374}{Publisher catalogue}.

\bibitem{Poonen}
Bjorn Poonen, Maximally complete fields, \emph{L'Enseignement
Math\'ematique} (2) \textbf{39} (1993), 87--106.
\href{https://math.mit.edu/~poonen/papers/amsval.pdf}{Author's manuscript},
especially the Hahn-series construction and Corollary~4.

\bibitem{Wang}
Wei Wang, \emph{A very short note on a problem of Godsil and Sun},
arXiv:2510.03816, 2025.
\href{https://arxiv.org/abs/2510.03816}{arXiv record}.
Cited as an explicit contemporary use of the classical real-closed-field
spectral theorem, not as the origin of that theorem.

\bibitem{Tao}
Terence Tao, \emph{254A, Notes 3a: Eigenvalues and sums of Hermitian
matrices}, January 12, 2010.
\href{https://terrytao.wordpress.com/2010/01/12/254a-notes-3a-eigenvalues-and-sums-of-hermitian-matrices/}{Author's lecture notes}.
Classical context for the variational principles; the field-general proofs
are supplied in the present article.

\bibitem{EY}
Carl Eckart and Gale Young, The approximation of one matrix by another of
lower rank, \emph{Psychometrika} \textbf{1} (1936), 211--218.
\href{https://doi.org/10.1007/BF02288367}{doi:10.1007/BF02288367}.

\bibitem{Penrose}
Roger Penrose, A generalized inverse for matrices,
\emph{Proceedings of the Cambridge Philosophical Society}
\textbf{51} (1955), 406--413.
\href{https://doi.org/10.1017/S0305004100030401}{doi:10.1017/S0305004100030401}.

\bibitem{DK}
Chandler Davis and William M. Kahan, The rotation of eigenvectors by a
perturbation. III, \emph{SIAM Journal on Numerical Analysis}
\textbf{7} (1970), 1--46.
\href{https://doi.org/10.1137/0707001}{doi:10.1137/0707001}.
Historical attribution; the Frobenius variant
used here is proved in \cref{thm:subspace}.

\bibitem{YWS}
Yi Yu, Tengyao Wang, and Richard J. Samworth,
\emph{A useful variant of the Davis--Kahan theorem for statisticians},
arXiv:1405.0680, 2014.
\href{https://arxiv.org/abs/1405.0680}{arXiv record}.
In particular the classical subspace perturbation framework and its
distinction between mixed-spectrum and unperturbed-spectrum gaps.

\bibitem{Adkins}
William A. Adkins, Normal matrices over hermitian discrete valuation rings,
\emph{Linear Algebra and its Applications} \textbf{157} (1991), 165--174.
\href{https://doi.org/10.1016/0024-3795(91)90111-9}{doi:10.1016/0024-3795(91)90111-9}.
A predecessor for the spectral theorem over formal and convergent
one-variable power series; conceded, not displaced.

\bibitem{KO}
Hans A. Keller and Herminia Ochsenius A., \emph{A spectral theorem for
matrices over fields of power series}, Annales math\'ematiques Blaise Pascal
\textbf{2}, no.~1 (1995), 169--179.
\href{https://ambp.centre-mersenne.org/item/AMBP_1995__2_1_169_0/}{Journal record}.
Orthogonal diagonalization over iterated real Laurent-series fields and one
infinite-rank Hahn field, by recursive block splitting.

\bibitem{PR}
Adam Parusi\'nski and Guillaume Rond, Multiparameter perturbation theory of
matrices and linear operators, \emph{Transactions of the American
Mathematical Society} \textbf{373} (2020), 2933--2948.
\href{https://doi.org/10.1090/tran/8061}{doi:10.1090/tran/8061};
\href{https://arxiv.org/abs/1807.04242}{arXiv:1807.04242}.
The formal separated-block lifting ingredient; the version used here is
proved independently in \cref{spec:thm:split}.

\bibitem{DLLZ}
Zhihong Dai, Haoran Liang, Jiaxin Lu, and Lihong Zhi,
\emph{An algorithm for diagonalizing matrices of formal power series},
arXiv:2602.08313v1, February 2026.
\href{https://arxiv.org/abs/2602.08313}{arXiv record}.
A preprint, not described here as a refereed publication; its ambient ring
is a multivariable formal power-series ring, not a Hahn field.

\bibitem{Neumann}
Bernhard H. Neumann, On ordered division rings, \emph{Transactions of the
American Mathematical Society} \textbf{66} (1949), 202--252.
\href{https://doi.org/10.1090/S0002-9947-1949-0032593-5}{doi:10.1090/S0002-9947-1949-0032593-5}.
The positive-support lemma, reproved as \cref{spec:lem:neumann}; the
bibliographic record rather than a full transcription was consulted.

\bibitem{Higman}
Graham Higman, Ordering by divisibility in abstract algebras,
\emph{Proceedings of the London Mathematical Society} (3) \textbf{2} (1952),
326--336.
\href{https://doi.org/10.1112/plms/s3-2.1.326}{doi:10.1112/plms/s3-2.1.326}.
The ordered-word special case needed here is proved as
\cref{spec:lem:words}.

\end{thebibliography}
\end{document}
